% This chapter is a faithful blueprint adaptation of the indicated part of
% Cao--Zhu's revised TeX source. Source-line sentinels preserve auditability.

\chapter{Ancient \texorpdfstring{$\kappa$}{kappa}-solutions}
\label{chap:cz-ancient-kappa-solutions}

\bigskip
Let us consider a solution of the Ricci flow on a compact
manifold. If the solution blows up in finite time (i.e., the
maximal solution exists only on a finite time interval), then as
we saw in Chapter 4 a sequence of rescalings of the solution
around the singularities converge to a solution which exists at
least on the time interval $(-\infty,T)$ for some finite number
$T$. Furthermore, by Perelman's no local collapsing theorem I
(Theorem 3.3.2), we see that the limit is $\kappa$-noncollapsed on
all scales for some positive constant $\kappa$. In addition, if
the dimension $n=3$ then the Hamilton-Ivey pinching estimate
implies that the limiting solution must have nonnegative curvature
operator.

We call a solution to the Ricci flow an \textbf{ancient
$\kappa$-solution}\index{ancient!$\kappa$-solution} if it is
complete (either compact or noncompact) and defined on an ancient
time interval $(-\infty,T)$ with $T>0$, has nonnegative curvature
operator and bounded curvature, and is $\kappa$-noncollapsed on
all scales for some positive constant $\kappa$.

In this chapter we study ancient $\kappa$-solutions of the Ricci
flow. We will obtain crucial curvature estimates of such solutions
and determine their structures in lower dimensional cases. In
particular, Sections 6.2-6.4 give a detailed exposition of
Perelman's work in section 11 of \cite{P1} and section 1 of
\cite{P2}. We also remak that the earlier work on ancient
solutions can be found in Hamilton \cite{Ha95F}.


\section{Preliminaries}

We first present a useful geometric property (cf. Proposition 2.2 of
\cite{CZ05F}) for complete noncompact Riemannian manifolds with
nonnegative sectional curvature.

Let $(M,g_{ij})$ be an $n$-dimensional complete Riemannian
manifold and let $\varepsilon$ be a positive constant. We call an
open subset $N\subset M$ an \textbf{$\varepsilon$-neck of radius
$r$}\index{$\varepsilon$-neck of radius $r$} if $(N,
r^{-2}g_{ij})$ is $\varepsilon$-close, in the
$C^{[\varepsilon^{-1}]}$ topology, to a standard neck
$\mathbb{S}^{n-1}\times \mathbb{I}$, where $\mathbb{S}^{n-1}$ is
the round $(n-1)$-sphere with scalar curvature 1 and $\mathbb{I}$
is an interval of length $2\varepsilon^{-1}$. The following result
is, to some extent, in similar spirit of Yau's volume lower bound
estimate \cite{Y76}.

%\vskip 0.2cm \noindent {\bf Proposition 6.1.1}
%CZ-SOURCE-LINE-13429
\begin{proposition}[lower radius bound for necks at infinity]
  \label{prop:cz-neck-radius-lower-bound-at-infinity}
  \dcref{Cao--Zhu Proposition 6.1.1}
  \group{Ancient-Solution Preliminaries}
  \source{cao-zhu}

There exists a positive constant $\varepsilon_0=\varepsilon_0(n)$
such that every complete noncompact $n$-dimensional Riemannian
manifold $(M, g_{ij})$ of nonnegative sectional curvature has a
positive constant $r_0$ such that any $\varepsilon$-neck of radius
$r$ on $(M, g_{ij})$ with $\varepsilon \leq \varepsilon_0$ must
have $r\geq r_0$.
\end{proposition}

%\vskip 0.1cm\noindent {\textbf{Proof}.}\
\begin{proof} The following argument is taken from \cite{CZ05F}.
We argue by contradiction. Suppose there exist a sequence of
positive constants $\varepsilon^\alpha \rightarrow 0$ and a
sequence of $n$-dimensional complete noncompact Riemannian
manifolds $(M^\alpha, g^{\alpha}_{ij})$ such that for each fixed
$\alpha$, there exists a sequence of $\varepsilon^\alpha$-necks
$N_{k}$ of radius at most $1/k$ in $M^\alpha$ with centers $P_{k}$
divergent to infinity.

Fix a point $P$ on the manifold $M^\alpha$ and connect each $P_k$
to $P$ by a minimizing geodesic $\gamma_k$. By passing to a
subsequence we may assume the angle $\theta_{kl}$ between geodesic
$\gamma_k$ and $\gamma_l$ at $P$ is very small and tends to zero
as $k, l\rightarrow+\infty$, and the length of $\gamma_{k+1}$ is
much bigger than the length of $\gamma_k$. Let us connect $P_k$ to
$P_l$ by a minimizing geodesic $\eta_{kl}$. For each fixed $l>k$,
let $\tilde{P}_k$ be a point on the geodesic $\gamma_l$ such that
the geodesic segment from $P$ to $\tilde{P}_k$ has the same length
as $\gamma_k$ and consider the triangle $\Delta PP_k\tilde{P}_k$
in $M^\alpha$ with vertices $P$, $P_k$ and $\tilde{P}_k$. By
comparing with the corresponding triangle in the Euclidean plane
$\mathbb{R}^2$ whose sides have the same corresponding lengths,
Toponogov's comparison theorem implies
$$
d(P_k,\tilde{P}_k)\le 2\sin\left(\frac{1}{2}\theta_{kl}\right)\cdot
d(P_k,P).
$$
Since $\theta_{kl}$ is very small, the distance from $P_k$ to the
geodesic $\gamma_l$ can be realized by a geodesic $\zeta_{kl}$
which connects $P_k$ to a point $P_k'$ on the interior of the
geodesic $\gamma_l$ and has length at most
$2\sin(\frac{1}{2}\theta_{kl})\cdot d(P_k,P).$ Clearly the angle
between $\zeta_{kl}$ and $\gamma_l$ at the intersection point
$P_k'$ is $\frac{\pi}{2}$. Consider $\alpha$ to be fixed and
sufficiently large. We claim that for large enough $k$, each
minimizing geodesic $\gamma_l$ with $l>k$, connecting $P$ to
$P_l$, goes through the neck $N_k$.

Suppose not; then the angle between $\gamma_k$ and $\zeta_{kl}$ at
$P_k$ is close to either zero or $\pi$ since $P_k$ is in the
center of an $\varepsilon^\alpha$-neck and $\alpha$ is
sufficiently large. If the angle between $\gamma_k$ and
$\zeta_{kl}$ at $P_k$ is close to zero, we consider the triangle
$\Delta PP_kP_k'$ in $M^\alpha$ with vertices $P$, $P_k$, and
$P_k'$. Note that the length between $P_k$ and $P_k'$ is much
smaller than the lengths from $P_k$ or $P_k'$ to $P$. By comparing
the angles of this triangle with those of the corresponding
triangle in the Euclidean plane with the same corresponding
lengths and using Toponogov's comparison theorem, we find that it
is impossible. Thus the angle between $\gamma_k$ and $\zeta_{kl}$
at $P_k$ is close to $\pi$. We now consider the triangle $\Delta
P_kP_k'P_l$ in $M^\alpha$ with the three sides $\zeta_{kl}$,
$\eta_{kl}$ and the geodesic segment from $P_k'$ to $P_l$ on
$\gamma_l$. We have seen that the angle of $\Delta P_k P_k'P_l$ at
$P_k$ is close to zero and the angle at $P_k'$ is $\frac{\pi}{2}$.
By comparing with corresponding triangle
$\bar{\Delta}\bar{P_k}\bar{P_k'}\bar{P_l}$ in the Euclidean plane
$\mathbb{R}^2$ whose sides have the same corresponding lengths,
Toponogov's comparison theorem implies
$$
\angle \bar{P_l}\bar{P_k}\bar{P_k'} +\angle
\bar{P_l}\bar{P_k'}\bar{P_k} \le\angle P_lP_kP_k'+\angle
P_lP_k'P_k<\frac{3}{4}\pi.
$$
This is impossible since the length between $\bar{P_k}$ and
$\bar{P_k'}$ is much smaller than the length from $\bar{P_l}$ to
either $\bar{P_k}$ or $\bar{P_k'}$. So we have proved each
$\gamma_l$ with $l>k$ passes through the neck $N_k$.

Hence by taking a limit, we get a geodesic ray $\gamma$ emanating
from $P$ which passes through all the necks $N_k$, $k = 1, 2,
\ldots,$  except a finite number of them. Throwing these finite
number of necks away, we may assume $\gamma$ passes through all
necks $N_k$, $k=1,2,\ldots.$ Denote the center sphere of $N_k$ by
$S_k$, and their intersection points with $\gamma$ by  $p_{k}\in
S_k\cap \gamma$, $k=1,2,\ldots.$


Take a sequence of points  $\gamma(m)$ with $m=1,2,\ldots.$ For
each fixed neck $N_k$, arbitrarily choose a point $q_{k}\in N_k$
near the center sphere $S_k$ and draw a geodesic segment
$\gamma^{km}$ from $q_{k}$ to $\gamma(m)$. Now we claim that for
any neck $N_l$ with $l>k$, $\gamma^{km}$ will pass through $N_l$
for all sufficiently large $m$.

We argue by contradiction. Let us place all the necks $N_i$
horizontally so that the geodesic $\gamma$ passes through each
$N_i$ from the left to the right. We observe that the geodesic
segment $\gamma^{km}$ must pass through the right half of $N_k$;
otherwise $\gamma^{km}$ cannot be minimal. Then for large enough
$m$, the distance from $p_{l}$ to the geodesic segment
$\gamma^{km}$ must be achieved by the distance from $p_l$ to some
interior point ${p_k}'$ of $\gamma^{km}$. Let us draw a minimal
geodesic $\eta$ from $p_{l}$ to the interior point ${p_k}'$ with
the angle at the intersection point ${p_k}'\in \eta\cap
\gamma^{km}$ to be $\frac{\pi}{2}.$ Suppose the claim is false.
Then the angle between $\eta$ and $\gamma$ at $p_{l}$ is close to
$0$ or $\pi$ since $\varepsilon^\alpha$ is small.

If the angle between $\eta$ and $\gamma$ at $p_{l}$ is close to
$0$, we consider the triangle ${\Delta} {p}_{l}{p_k}'{\gamma}(m)$
and construct a comparison triangle $ \bar{\Delta}
\bar{p}_{l}\bar{p_k}'\bar{\gamma}(m)$ in the plane with the same
corresponding length. Then by Toponogov's comparison theorem, we
see the sum of the inner angles of the comparison triangle $
\bar{\Delta} \bar{p}_{l}\bar{p_k}'\bar{\gamma}(m)$ is less than
$3\pi/4$, which is impossible.

If the angle between $\eta$ and $\gamma$ at $p_{l}$ is close to
$\pi$, by drawing a minimal geodesic $\xi$ from $q_k$ to $p_{l}$,
we see that $\xi$ must pass through the right half of $N_k$ and
the left half of $N_l$; otherwise $\xi$ cannot be minimal. Thus
the three inner angles of the triangle $\Delta p_{l}{p_k}'q_k$ are
almost $0,\pi/2$, and $0$ respectively. This is also impossible by
the Toponogov comparison theorem.

Hence we have proved that the geodesic segment $\gamma^{km}$
passes through $N_l$ for $m$ large enough.

Consider the triangle $\Delta p_{k}q_k\gamma(m)$ with two long
sides $\overline{p_{k}\gamma(m)}(\subset\gamma)$ and
$\overline{q_{k}\gamma(m)}(= \gamma^{km})$. For any $s>0$, choose
points ${\tilde{p}_{k}}$ on $\overline{p_{k}\gamma(m)}$ and
${\tilde{q}_{k}}$ on $\overline{q_{k}\gamma(m)}$ with
$d(p_{k},{\tilde{p}_{k}})=d(q_{k},{\tilde{q}_{k}})=s$. By
Toponogov's comparison theorem, we have {\small
\begin{align*}
&  \left(\frac{d({\tilde{p}_{k}},{\tilde{q}_{k}})}{d(p_{k},q_{k})}\right)^{2}\\
&  =\frac{d({\tilde{p}_{k}},\gamma(m))^{2}
+d({\tilde{q}_{k}},\gamma(m))^{2}-2d({\tilde{p}_{k}},
\gamma(m))d({\tilde{q}_{k}}, \gamma(m))\cos\bar{\measuredangle}
({\tilde{p}_{k}}\gamma(m){\tilde{q}_{k}})}
{d({p_{k}},\gamma(m))^{2}+d({q_{k}},\gamma(m))^{2}-2d({p_{k}},
\gamma(m))d({q_{k}},\gamma(m))\cos\bar{\measuredangle}
({p_{k}}\gamma(m){q_{k}})}\\[1mm]
&  \geq \frac{d({\tilde{p}_{k}},\gamma(m))^{2}+d({\tilde{q}_{k}},
\gamma(m))^{2}-2d({\tilde{p}_{k}},\gamma(m))d({\tilde{q}_{k}},
\gamma(m))\cos\bar{\measuredangle}
({\tilde{p}_{k}}\gamma(m){\tilde{q}_{k}})}
{d({p_{k}},\gamma(m))^{2}+d({q_{k}},\gamma(m))^{2}
-2d({p_{k}},\gamma(m))d({q_{k}},
\gamma(m))\cos\bar{\measuredangle}
({\tilde{p}_{k}}\gamma(m){\tilde{q}_{k}})}\\[1mm]
&  = \frac{(d({\tilde{p}_{k}},\gamma(m)) -
d({\tilde{q}_{k}},\gamma(m)))^{2}+2d({\tilde{p}_{k}},
\gamma(m))d({\tilde{q}_{k}}, \gamma(m))(1 -
\cos\bar{\measuredangle}
({\tilde{p}_{k}}\gamma(m){\tilde{q}_{k}}))}
{(d({\tilde{p}_{k}},\gamma(m)) -
d({\tilde{q}_{k}},\gamma(m)))^{2}+2d({{p}_{k}},\gamma(m))d({{q}_{k}},
\gamma(m))(1 - \cos\bar{\measuredangle}
({\tilde{p}_{k}}\gamma(m){\tilde{q}_{k}}))}\\[1mm]
&  \geq \frac{d({\tilde{p}_{k}},\gamma(m))d({\tilde{q}_{k}},
\gamma(m))} {d({p_{k}},\gamma(m))d({q_{k}}, \gamma(m))}\\[1mm]
&   \rightarrow 1
\end{align*}
} as $m\rightarrow\infty$, where $\bar{\measuredangle}
({{p}_{k}}\gamma(m){{q}_{k}})$ and $\bar{\measuredangle}
({\tilde{p}_{k}}\gamma(m){\tilde{q}_{k}})$ are the corresponding
angles of the comparison triangles.

Letting $m\rightarrow \infty$, we see that $\gamma^{km}$ has a
convergent subsequence whose limit $\gamma^{k}$ is a geodesic ray
passing through all $N_l$ with $l > k$. Let us denote by
$p_{j}=\gamma(t_j), j=1, 2, \ldots$. From the above computation,
we deduce that
$$
d(p_{k},q_{k})\leq d(\gamma(t_k+s),\gamma^{k}(s))
$$
for all $s>0$.

Let $\varphi(x)=\lim_{t\rightarrow+\infty}(t-d(x,\gamma(t)))$ be
the Busemann function constructed from the ray $\gamma$. Note that
the level set $\varphi^{-1}(\varphi(p_{j}))\cap N_j$ is close to
the center sphere $S_j$ for any $j = 1, 2, \ldots$. Now let $q_k$
be any fixed point in $\varphi^{-1}(\varphi(p_{k}))\cap N_k$. By
the definition of Busemann function $\varphi$ associated to the
ray $\gamma$, we see that
$\varphi(\gamma^{k}(s_1))-\varphi(\gamma^{k}(s_2))=s_1-s_2$ for
any $s_1$, $s_2\geq0$. Consequently, for each $l>k$, by choosing
$s=t_l-t_k$, we see $\gamma^{k}(t_l-t_k)\in
\varphi^{-1}(\varphi(p_{l}))\cap N_l.$ Since
$\gamma(t_k+t_l-t_k)=p_{l}$, it follows that
$$
d(p_{k},q_{k})\leq d(p_{l},\gamma^{k}(s)).
$$
with $s =t_l-t_k >0$. This implies that the diameter of
$\varphi^{-1}(\varphi(p_{k}))\cap N_k$ is not greater than the
diameter of $\varphi^{-1}(\varphi(p_{l}))\cap N_l$ for any $l>k$,
which is a contradiction for $l$ much larger than $k$.

Therefore we have proved the proposition.
\end{proof}
%$$\eqno \#$$\vskip 0.3cm

In \cite{Ha95F}, Hamilton discovered an important repulsion
principle (cf. Theorem 21.4 of \cite{Ha95F}) about the influence
of a bump of strictly positive curvature in a complete noncompact
manifold of nonnegative sectional curvature. Namely minimal
geodesic paths that go past the bump have to avoid it. As a
consequence he obtained a finite bump theorem (cf. Theorem 21.5 of
\cite{Ha95F}) that gives a bound on the number of bumps of
curvature.

Let $M$ be a complete noncompact Riemannian manifold with
nonnegative sectional curvature $K\ge 0$. A geodesic ball $B(p,r)$
of radius $r$ centered at a point $p\in M$ is called a {\bf
curvature $\beta$-bump}\index{curvature $\beta$-bump} if sectional
curvature $K\ge \beta/r^2$ at all points in the ball. The ball
$B(p,r)$ is called {\bf $\lambda$-remote}\index{$\lambda$-remote}
from an origin $O$ if $d(p, O)\ge \lambda r$.

\medskip
{\bf Finite Bump Theorem}\index{finite bump theorem} (Hamilton
\cite{Ha95F}){\bf .} \emph{ For every $\beta >0$ there exists
$\lambda<\infty$ such that in any complete manifold of nonnegative
sectional curvature there are at most a finite number of disjoint
balls which are $\lambda$-remote curvature $\beta$-bumps.}

\medskip
This finite bump theorem played an important role in Hamilton's
study of the behavior of singularity models at infinity and in the
dimension reduction argument he developed for the Ricci flow (cf.
Section 22 of \cite{Ha95F}, see also \cite{CTZ04} for application
to the K\"ahler-Ricci flow and uniformization problem in complex
dimension two). A special consequence of the finite bump theorem
is that if we have a complete noncompact solution to the Ricci
flow on an ancient time interval $-\infty<t<T$ with $T>0$
satisfying certain local injectivity radius bound, with curvature
bounded at each time and with asymptotic scalar curvature ratio
$A=\limsup Rs^2=\infty$, then we can find a sequence of points
$p_j$ going to $\infty$ (as in the following Lemma 6.1.3) such
that a cover of the limit of dilations around these points at time
$t=0$ splits as a product with a flat factor. The following
result, somewhat known in Alexandrov space theory (e.g., Perelman
wrote in \cite{P1} (page 29, line 3) ``this follows from a
standard argument based on the Aleksandrov-Toponogov concavity"
and the essentially same statement also appeared in \cite{KL},
\cite{CZ05F} and \cite{CLN}), is in similar spirit as Hamilton's
finite bumps theorem and its consequence. The advantage is that we
will get in the limit of dilations a product of the line with a
lower dimensional manifold, instead of a quotient of such a
product.

%\vskip 0.2cm \noindent {\bf Proposition 6.1.2}
%CZ-SOURCE-LINE-13684
\begin{proposition}[line splitting in pointed limits at infinity]
  \label{prop:cz-pointed-limit-at-infinity-splits-line}
  \dcref{Cao--Zhu Proposition 6.1.2}
  \group{Ancient-Solution Preliminaries}
  \source{cao-zhu}

Suppose $(M,g_{ij})$ is a complete $n$-dimensional Riemannian
manifold with nonnegative sectional curvature. Let $P\in M$ be
fixed, and $P_k\in M$ a sequence of points and $\lambda_k$ a
sequence of positive numbers with $d(P,P_k)\rightarrow +\infty$
and $\lambda_kd(P,P_k)\rightarrow +\infty$. Suppose also that the
marked manifolds $(M,\lambda_k^2g_{ij},P_k)$ converge in the
$C^{\infty}_{loc}$ topology to a Riemannian manifold
$\widetilde{M}$. Then the limit $\widetilde{M}$ splits
isometrically as the metric product of the form $\mathbb{R}\times
N$, where $N$ is a Riemannian manifold with nonnegative sectional
curvature.
\end{proposition}

%\vskip 0.1cm\noindent {\textbf{Proof}.}\
\begin{proof}
We now follow an argument given in \cite{CZ05F}. Let us denote by
$|OQ|=d(O,Q)$ the distance between two points $O,Q\in M$. Without
loss of generality, we may assume that for each $k$, \begin{equation}
1+2|PP_k|\leq|PP_{k+1}|.
%\eqno (6.1.1)
\end{equation} Draw a minimal geodesic $\gamma_k$ from $P$ to $P_k$ and a
minimal geodesic $\sigma_k$ from $P_k$ to $P_{k+1}$, both
parametrized by arclength. We may further assume \begin{equation}
\theta_k=|\measuredangle(\dot{\gamma}_k(0),\dot{\gamma}_{k+1}(0))
|<\frac{1}{k}.
%\eqno (6.1.2)
\end{equation}

By assumption, the sequence $(M,\lambda^2_kg_{ij},P_k)$ converges
(in the $C^{\infty}_{loc}$ topology) to a Riemannian manifold
$(\widetilde{M},\widetilde{g}_{ij},\widetilde{P})$ with
nonnegative sectional curvature. By a further choice of
subsequences, we may also assume $\gamma_k$ and $\sigma_k$
converge to geodesic rays $\widetilde{\gamma}$ and
$\widetilde{\sigma}$ starting at $\widetilde{P}$ respectively. We
will prove that $\tilde{\gamma}\cup\tilde{\sigma}$ forms a line in
$\widetilde{M}$, and then by the Toponogov splitting theorem
\cite{Mi} the limit $\widetilde{M}$ must be splitted as
$\mathbb{R}\times N$.

We argue by contradiction. Suppose
$\widetilde{\gamma}\cup\widetilde{\sigma}$ is not a line; then for
each $k$, there exist two points $A_k\in \gamma_k$ and $B_k\in
\sigma_k$ such that as $k\rightarrow +\infty$, \begin{equation} \left\{
\begin{array}{llll}
\lambda_kd(P_k,A_k)\rightarrow A>0,\\[2mm]
\lambda_kd(P_k,B_k)\rightarrow B>0,\\[2mm]
\lambda_kd(A_k,B_k)\rightarrow C>0,\\[2mm]
\mbox{but  }A+B>C.\\
\end{array}\right. %\eqno(6.1.3)
\end{equation}

\begin{center}
\setlength{\unitlength}{2mm}
\begin{picture}(60,20)
\linethickness{1pt} \put(0,0){\line(6,1){60}}

\put(0,0){\line(2,1){25}} \thicklines

\qbezier(25,12.5)(40,11.5)(60,10)

\put(20,10){\line(6,1){12}}

\put(0,-2){\makebox(2,1)[c]{$P$}}

\put(25,15){\makebox(2,1)[c]{$P_k$}}

\put(60,8){\makebox(2,1)[c]{$P_{k+1}$}}

\put(20,8){\makebox(2,1)[c]{$A_k$}}

\put(25,9){\makebox(2,1)[c]{$\delta_k$}}

\put(32,10){\makebox(2,1)[c]{$B_k$}}

\put(42,12){\makebox(2,1)[c]{$\sigma_k$}}

\put(12,7){\makebox(2,1)[c]{$\gamma_k$}}

\end{picture}

\end{center}

\bigskip\bigskip
Now draw a minimal geodesic $\delta_k$ from $A_k$ to $B_k$.
Consider comparison triangles
$\bar{\triangle}{\bar{P}_k\bar{P}\bar{P}_{k+1}}$ and
$\bar{\triangle}{\bar{P}_k\bar{A}_k\bar{B}_k}$ in $\mathbb{R}^2$
with
$$
|\bar{P}_k\bar{P}|=|P_kP|, |\bar{P}_k\bar{P}_{k+1}|=|P_kP_{k+1}|,
|\bar{P}\bar{P}_{k+1}|=|PP_{k+1}|,
$$
$$\mbox{ and }\;|\bar{P}_k\bar{A}_k|
=|P_kA_k|, |\bar{P}_k\bar{B}_k|=|P_kB_k|,
|\bar{A}_k\bar{B}_k|=|A_kB_k|.
$$
By Toponogov's comparison theorem \cite{BGP}, we have \begin{equation}
\measuredangle\bar{A}_k\bar{P}_k\bar{B}_k\geq
\measuredangle \bar{P} \bar{{P}_k}\bar{P}_{k+1}. %\eqno (6.1.4)
\end{equation} On the other hand, by (6.1.2) and using Toponogov's comparison
theorem again, we have \begin{equation}
\measuredangle\bar{P}_k\bar{P}\bar{P}_{k+1} \leq \measuredangle
P_kPP_{k+1}<\frac{1}{k},
%\eqno(6.1.5)
\end{equation} and since $|\bar{P}_k\bar{P}_{k+1}|>|\bar{P}\bar{P}_k|$ by
(6.1.1), we further have \begin{equation}
\measuredangle\bar{P}_k\bar{P}_{k+1}\bar{P} \leq
\measuredangle\bar{P}_k\bar{P}\bar{P}_{k+1}<\frac{1}{k}.
%\eqno(6.1.6)
\end{equation} Thus the above inequalities (6.1.4)-(6.1.6) imply that
$$
\measuredangle\bar{A}_k\bar{P}_k\bar{B}_k>\pi-\frac{2}{k}.
$$
Hence \begin{equation} |\bar{A}_k\bar{B}_k|^2\geq
|\bar{A}_k\bar{P}_k|^2+|\bar{P}_k\bar{B}_k|^2
-2|\bar{A}_k\bar{P}_k|\cdot|\bar{P}_k\bar{B}_k|\cos\left(\pi-\frac{2}{k}\right).
%\eqno (6.1.7)
\end{equation}

Multiplying the above inequality by $\lambda^2_k$ and letting
$k\rightarrow+\infty$, we get $$C\geq A+B$$ which contradicts
(6.1.3).

Therefore we have proved the proposition.
%$$\eqno \#$$\vskip 0.3cm
\end{proof}

Let $M$ be an $n$-dimensional complete noncompact Riemannian
manifold. Pick an origin $O\in M$. Let $s$ be the geodesic
distance to the origin $O$ of $M$, and $R$ the scalar curvature.
Recall that in Chapter 4 we have defined the {\bf asymptotic
scalar curvature ratio}\index{asymptotic scalar curvature ratio}
$$
A=\limsup_{s\rightarrow+\infty}Rs^2.
$$

We now state a useful lemma of Hamilton (cf. Lemma 22.2 of
\cite{Ha95F}) about picking local (almost) maximum curvature
points at infinity.

%\vskip 0.2cm \noindent{\bf Lemma 6.1.3}\ \
%CZ-SOURCE-LINE-13827
\begin{lemma} [point selection at infinite asymptotic curvature ratio]
  \label{lem:cz-infinite-ascr-point-selection}
  \dcref{Cao--Zhu Lemma 6.1.3}
  \group{Ancient-Solution Preliminaries}
  \source{cao-zhu}

 Given a complete noncompact Riemannian manifold with
bounded curvature and with asymptotic scalar curvature ratio
$$
A=\limsup_{s\rightarrow+\infty}Rs^2=+\infty,
$$
we can find a sequence of points $x_j$ divergent to infinity, a
sequence of radii $r_j$ and a sequence of positive numbers
$\delta_j\rightarrow0$ such that
\begin{itemize}
\item[(i)] $R(x)\leq(1+\delta_j)R(x_j)$ for all $x$ in the ball
$B(x_j,r_j)$ of radius $r_j$ around $x_j$, \item[(ii)]
$r^2_jR(x_j)\rightarrow+\infty$, \item[(iii)]
$\lambda_j=d(x_j,O)/r_j\rightarrow+\infty$, \item[(iv)]
the balls $B(x_j,r_j)$ are disjoint, \\
where $d(x_j,O)$ is the distance of $x_j$ from the origin $O$.
\end{itemize}
\end{lemma}

%\vskip 0.1cm\noindent {\textbf{Proof}.}\
\begin{proof} The proof is essentially from Hamilton \cite{Ha95F}.
Pick a sequence of positive numbers $\epsilon_j\rightarrow0$, then
choose $A_j\rightarrow+\infty$ so that
$A_j\epsilon^2_j\rightarrow+\infty$. Let $\sigma_j$ be the largest
number such that
$$
\sup\{ R(x)d(x,O)^2\ |\ d(x,O)\leq\sigma_j\}\leq A_j.
$$
Then there exists some $y_j\in M$ such that
$$
R(y_j)d(y_j,O)^2=A_j\ \ \ \mbox{and}\ \ \ d(y_j,O)=\sigma_j.
$$
Now pick $x_j\in M$ so that $d(x_j,O)\geq\sigma_j$ and
$$
R(x_j)\geq\frac{1}{1+\epsilon_j}\sup\{ R(x)\ |\
d(x,O)\geq\sigma_j\}.
$$
Finally pick $r_j=\epsilon_j\sigma_j$. We check the properties
(i)-(iv) as follows.

\smallskip
(i)\ \ \ If $x\in B(x_j,r_j)\cap\{d(\cdot,O)\geq\sigma_j\}$, we
have
$$
R(x)\leq(1+\epsilon_j)R(x_j)
$$
by the choice of the point $x_j$; while if $x\in
B(x_j,r_j)\cap\{d(\cdot,O)\leq\sigma_j\}$, we have
\begin{displaymath}
\begin{split}
   R(x)&\leq A_j/d(x,O)^2\\[1mm]
       &\leq \frac{1}{(1-\epsilon_j)^2}(A_j/\sigma_j^2)\\[1mm]
       &=  \frac{1}{(1-\epsilon_j)^2}R(y_j)\\[1mm]
       &\leq \frac{(1+\epsilon_j)}{(1-\epsilon_j)^2}R(x_j),
\end{split}
\end{displaymath}
since $d(x,O)\geq
d(x_j,O)-d(x,x_j)\geq\sigma_j-r_j=(1-\epsilon_j)\sigma_j$. Thus we
have obtained
$$
R(x)\leq(1+\delta_j)R(x_j),\ \ \ \forall\ x\in B(x_j,r_j),
$$
where
$\delta_j=\frac{(1+\epsilon_j)}{(1-\epsilon_j)^2}-1\rightarrow0$
as $j\rightarrow+\infty$.

\smallskip
(ii)\ \ \ By the choices of $r_j$, $x_j$ and $y_j$, we have
\begin{displaymath}
\begin{split}
r^2_jR(x_j)&= \epsilon_j^2\sigma_j^2R(x_j)\\[1mm]
      &\geq \epsilon_j^2\sigma_j^2
\left[\frac{1}{1+\epsilon_j}R(y_j)\right]\\[1mm]
      &= \frac{\epsilon_j^2}{1+\epsilon_j}A_j\rightarrow+\infty,\
      \ \ \mbox{as }\; j\rightarrow+\infty.
\end{split}
\end{displaymath}

\smallskip
(iii)\ \ \ Since $d(x_j,O)\geq\sigma_j={r_j}/{\epsilon_j}$, it
follows that $\lambda_j={d(x_j,O)}/{r_j}\rightarrow+\infty$ as
$j\rightarrow+\infty$.

\smallskip
(iv)\ \ \ For any $x\in B(x_j,r_j)$, the distance from the origin
\begin{displaymath}
\begin{split}
   d(x,O)&\geq d(x_j,O)-d(x,x_j)\\[1mm]
         &\geq \sigma_j-r_j\\[1mm]
         &= (1-\epsilon_j)\sigma_j\rightarrow+\infty,\ \ \mbox{as }\;
         j\rightarrow+\infty.
\end{split}
\end{displaymath}
Thus any fixed compact set does not meet the balls $B(x_j,r_j)$
for large enough $j$. If we pass to a subsequence, the balls will
all avoid each other.
%$$\eqno \#$$ \vskip 0.3cm
\end{proof}

The above point picking lemma of Hamilton, as written down in
Lemma 22.2 of \cite{Ha95F}, requires the curvature of the manifold
to be bounded. When the manifold has unbounded curvature, we will
appeal to the following simple lemma.

%\vskip 0.2cm \noindent{\bf Lemma 6.1.4}\ \ \emph{
%CZ-SOURCE-LINE-13932
\begin{lemma}[point selection for unbounded curvature]
  \label{lem:cz-unbounded-curvature-point-selection}
  \dcref{Cao--Zhu Lemma 6.1.4}
  \group{Ancient-Solution Preliminaries}
  \source{cao-zhu}

Given a complete noncompact Riemannian manifold with unbounded
curvature, we can find a sequence of points $x_j$ divergent to
infinity such that for each positive integer $j$, we have
$|Rm(x_{j})|\geq j$, and
$$
|Rm(x)|\leq4|Rm(x_{j})|
$$
for $x\in B(x_j,\frac{j}{\sqrt{|Rm(x_{j})|}})$.
\end{lemma}

%\vskip 0.1cm\noindent {\textbf{Proof.}}\
\begin{proof}
Each $x_j$ can be constructed as a limit of a finite sequence
$\{y_i\}$, defined as follows. Let $y_0$ be any fixed point with
$|Rm(y_{0})| \geq j$. Inductively, if $y_i$ cannot be taken as
$x_j$, then there is a $y_{i+1}$ such that
\begin{equation*}
 \left\{
\begin{aligned}
|Rm(y_{i+1})|&  > 4|Rm(y_{i})|,\\
d(y_i,y_{i+1}) &  \leqslant \frac{j}{\sqrt{|Rm(y_{i})|}}.
\end{aligned}
\right.
\end{equation*}
Thus we have
$$
|Rm(y_{i})|> 4^i|Rm(y_{0})|\geq 4^ij,$$
$$d(y_i,y_{0})\leq
j\sum^i_{k=1}\frac{1}{\sqrt{4^{k-1}j}}<2\sqrt{j}.$$ Since the
manifold is smooth, the sequence $\{y_i\}$ must be finite. The
last element fits.
%$$\eqno \#$$ \vskip 0.5cm
\end{proof}


\section{Asymptotic Shrinking Solitons}
The main purpose of this section is to prove a result of Perelman
(cf. Proposition 11.2 of \cite{P1}) on the asymptotic shapes of
ancient $\kappa$-solutions as time $t \rightarrow -\infty$.


We begin with the study of the asymptotic behavior of an ancient
$\kappa$-solution $g_{ij}(x,t)$, on $M\times (-\infty,T)$ with
$T>0$, to the Ricci flow as $t\rightarrow-\infty$.

Pick an arbitrary point $(p,t_0)\in M\times(-\infty,0]$ and recall
from Chapter 3 that
$$
\tau=t_0-t,\ \ \ \mbox{for }\; t<t_0,
$$
\begin{multline*}
l(q,\tau)=\frac{1}{2\sqrt{\tau}}\inf \bigg\{
\int_0^\tau\sqrt{s}\Big(R(\gamma(s),t_0-s) \\
+|\dot{\gamma}(s)|^2_{g_{ij}(t_0-s)}\Big)ds \left|
\begin{array}{c}
\gamma:[0,\tau]\rightarrow M\ \mbox{with}\\[1mm]
\gamma(0)=p,\ \gamma(\tau)=q\
\end{array}\bigg\}\right.
\end{multline*}
and
$$
\tilde{V}(\tau)=\int_M(4\pi\tau)^{-\frac{n}{2}}
\exp(-l(q,\tau))dV_{t_0-\tau}(q).
$$
We first observe that Corollary 3.2.6 also holds for the general
complete manifold $M$. Indeed, since the scalar curvature is
nonnegative, the function $\bar{L}(\cdot,\tau)=4\tau
l(\cdot,\tau)$ achieves its minimum on $M$ for each fixed
$\tau>0$. Thus the same argument in the proof of Corollary 3.2.6
shows there exists $q=q(\tau)$ such that \begin{equation}
l(q(\tau),\tau)\leq\frac{n}{2} %\eqno(6.2.1)
\end{equation} for each $\tau>0$.

Recall from (3.2.11)-(3.2.13), the Li-Yau-Perelman distance $l$
satisfies the following
\begin{align}
\frac{\partial}{\partial\tau}l
&=-\frac{l}{\tau}+R+\frac{1}{2\tau^{3/2}}K, \\  %\eqno(6.2.2)
|\nabla l|^2&=-R+\frac{l}{\tau}-\frac{1}{\tau^{3/2}}K, \\  %\eqno(6.2.3)
\Delta l& \leq-R+\frac{n}{2\tau}-\frac{1}{2\tau^{3/2}}K,  %\eqno(6.2.4)
\end{align}
and the equality in (6.2.4) holds everywhere if and only if we are
on a gradient shrinking soliton. Here $K=\int_0^\tau
s^{3/2}Q(X)ds,$ $Q(X)$ is the trace Li-Yau-Hamilton quadratic
given by
$$
Q(X)=-R_{\tau} -\frac{R}{\tau}-2 \langle\nabla
R,X\rangle+2\Ric(X,X)
$$
and $X$ is the tangential (velocity) vector field of an
$\mathcal{L}$-shortest curve $\gamma: [0,\tau]\rightarrow M$
connecting $p$ to $q$.

By applying the trace Li-Yau-Hamilton inequality (Corollary 2.5.5)
to the ancient $\kappa$-solution, we have
\begin{align*}
Q(X)&  = -R_{\tau} -\frac{R}{\tau}-2 \langle\nabla
R,X\rangle+2\Ric(X,X) \\
&  \geq -\frac{R}{\tau}
\end{align*}
and hence
\begin{align*}
K&  = \int_0^{\tau}s^{3/2}Q(X)ds   \\
&  \geq -\int_0^{\tau}\sqrt{s}Rds \\
&  \geq -L(q,\tau).
\end{align*}
Thus by (6.2.3) we get \begin{equation} |\nabla l|^2+R\leq \frac{3l}{\tau}.
%\eqno(6.2.5)
\end{equation} %\vskip 0.3cm

\medskip
We are now state and prove the following

%\vskip 0.2cm\noindent{\bf Theorem 6.2.1}
%CZ-SOURCE-LINE-14047
\begin{theorem}[asymptotic shrinking soliton of an ancient kappa-solution]
  \label{thm:cz-ancient-kappa-asymptotic-shrinker}
  \dcref{Cao--Zhu Theorem 6.2.1}
  \group{Asymptotic Shrinking Solitons}
  \source{cao-zhu}

Let $g_{ij}(\cdot,t), -\infty<t<T$ with some $T>0$, be a nonflat
ancient $\kappa$-solution for some $\kappa >0$. Then there exist a
sequence of points $q_k$ and a sequence of times $t_k \rightarrow
-\infty$ such that the scalings of $g_{ij}(\cdot,t)$ around $q_k$
with factor $|t_k|^{-1}$ and with the times $t_k$ shifting to the
new time zero converge to a nonflat gradient shrinking soliton in
$C_{\rm loc}^\infty$ topology.
\end{theorem}

%\vskip 0.1cm \noindent {\bf Proof.}
\begin{proof}
  \uses{lem:cz-reduced-distance-curvature-parabolic-annulus} The proof basically follows the argument of Perelman
(11.2 of \cite{P1}).  Clearly, we may assume that the nonflat
ancient $\kappa$-solution is not a gradient shrinking soliton. For
the arbitrarily fixed $(p,t_0)$, let $q(\tau) (\tau=t_0 -t)$ be
chosen as in (6.2.1) with $l(q(\tau),\tau) \leq \frac{n}{2}$. We
only need to show that the scalings of $g_{ij}(\cdot,t)$ around
$q(\tau)$ with factor $\tau^{-1}$ converge along a subsequence of
$\tau \rightarrow +\infty$ to a nonflat gradient shrinking soliton
in the $C^{\infty}_{\rm loc}$ topology.

We first claim that for any $A \geq 1$, one can find
$B=B(A)<+\infty$ such that for every $\bar{\tau}>1$ there holds
\begin{equation} l(q,\tau)\leq B \; \mbox{ and }\;  \tau R(q,t_0-\tau)\leq B,
%\eqno(6.2.6)
\end{equation} whenever $\frac{1}{2}\bar{\tau}\leq \tau \leq A\bar{\tau}$ and
$d_{t_0-\frac{\bar{\tau}}{2}}^2 (q,q(\frac{\bar{\tau}}{2}))\leq
A\bar{\tau}.$

Indeed, by using (6.2.5) at $\tau=\frac{\bar{\tau}}{2}$, we have
\begin{align}
\sqrt{l(q,\frac{\bar{\tau}}{2})} &  \leq \sqrt{\frac{n}{2}}+
\sup\{|\nabla \sqrt{l}|\}\cdot d_{t_0-\frac{\bar{\tau}}{2}}
\left(q,q\left(\frac{\bar{\tau}}{2}\right)\right) \\
&  \leq \sqrt{\frac{n}{2}}+ \sqrt{\frac{3}{2\bar{\tau}}}\cdot
\sqrt{A\bar{\tau}} \nonumber\\
&  = \sqrt{\frac{n}{2}}+\sqrt{\frac{3A}{2}},\nonumber
\end{align}  %\eqno(6.2.7)
and
\begin{align}
R\left(q,t_0-\frac{\bar{\tau}}{2}\right)&  \leq
\frac{3l(q,\frac{\bar{\tau}}{2})}{(\frac{\bar{\tau}}{2})}\\
&  \leq
\frac{6}{\bar{\tau}}\left(\sqrt{\frac{n}{2}}+\sqrt{\frac{3A}{2}}\right)^2,\nonumber
\end{align}  %\eqno (6.2.8)
for $q\in
B_{t_0-\frac{\bar{\tau}}{2}}(q(\frac{\bar{\tau}}{2}),\sqrt{A\bar{\tau}})$.
Recall that the Li-Yau-Hamilton inequality implies that the scalar
curvature of the ancient solution is pointwise nondecreasing in
time. Thus we know from (6.2.8) that \begin{equation} \tau R(q,t_0-\tau)\leq
6A\left(\sqrt{\frac{n}{2}}+\sqrt{\frac{3A}{2}}\right)^2 %\eqno(6.2.9)
\end{equation} whenever $\frac{1}{2}\bar{\tau}\leq \tau \leq A\bar{\tau}$ and
$d_{t_0-\frac{\bar{\tau}}{2}}^2 (q,q(\frac{\bar{\tau}}{2}))\leq
A\bar{\tau}.$

By (6.2.2) and (6.2.3) we have
$$
\frac{\partial l}{\partial \tau}+\frac{1}{2}|\nabla
l|^2=-\frac{l}{2\tau}+\frac{R}{2}.
$$
This together with (6.2.9) implies that
\begin{align*}
\frac{\partial l}{\partial \tau} &\leq
-\frac{l}{2\tau}+\frac{3A}{\tau}
\left(\sqrt{\frac{n}{2}}+\sqrt{\frac{3A}{2}}\right)^2 \\
\intertext{i.e.,} \frac{\partial}{\partial\tau}(\sqrt{\tau}l)
&\leq\frac{3A}{\sqrt{\tau}}\left(\sqrt{\frac{n}{2}}+\sqrt{\frac{3A}{2}}\right)^2
\end{align*}
whenever $\frac{1}{2}\bar{\tau}\leq \tau \leq A\bar{\tau}$ and
$d_{t_0-\frac{\bar{\tau}}{2}}^2
\left(q,q\left(\frac{\bar{\tau}}{2}\right)\right)\leq A\bar{\tau}.$ Hence by
integrating this differential inequality, we obtain
$$
\sqrt{\tau}l(q,\tau)-\sqrt{\frac{\bar{\tau}}{2}}l\left(q,\frac{\bar{\tau}}{2}\right)
\leq 6A\left(\sqrt{\frac{n}{2}}+\sqrt{\frac{3A}{2}}\right)^2\sqrt{\tau}
$$
and then by (6.2.7),
\begin{align}
l(q,\tau)&  \leq
l\left(q,\frac{\bar{\tau}}{2}\right)+6A\left(\sqrt{\frac{n}{2}}+\sqrt{\frac{3A}{2}}\right)^2\\
&  \leq 7A\left(\sqrt{\frac{n}{2}}+\sqrt{\frac{3A}{2}}\right)^2 \nonumber
\end{align}  %\eqno (6.2.10)
whenever $\frac{1}{2}\bar{\tau}\leq \tau \leq A\bar{\tau}$ and
$d_{t_0-\frac{\bar{\tau}}{2}}^2 (q,q(\frac{\bar{\tau}}{2}))\leq
A\bar{\tau}.$  So we have proved claim (6.2.6).

Recall that $g_{ij}(\tau)=g_{ij}(\cdot,t_0-\tau) $ satisfies
$(g_{ij})_\tau=2R_{ij}$. Let us take the scaling of the ancient
$\kappa$-solution around $q(\frac{\bar{\tau}}{2})$ with factor
$(\frac{\bar{\tau}}{2})^{-1}$, i.e.,
$$
\tilde{g}_{ij}(s)=\frac{2}{\bar{\tau}}g_{ij}
\left(\cdot,t_0-s\frac{\bar{\tau}}{2}\right)
$$
where $s\in [0,+\infty)$. Claim (6.2.6) says that for all
$s\in[1,2A]$ and all $q$ such that
dist$^2_{\tilde{g}_{ij}(1)}(q,q(\frac{\bar{\tau}}{2}))\leq A,$ we
have $\tilde{R}(q,s)=\frac{\bar
\tau}{2}R(q,t_0-s\bar{\frac{\tau}{2}})\leq B$. Now taking into
account the $\kappa$-noncollapsing assumption and Theorem 4.2.2,
we can use Hamilton's compactness theorem (Theorem 4.1.5) to
obtain a sequence $\bar{\tau}_k\rightarrow +\infty$ such that the
marked evolving manifolds
$(M,\tilde{g}_{ij}^{(k)}(s),q(\frac{\bar{\tau}_k}{2})),$ with
$\tilde{g}_{ij}^{(k)}(s)=\frac{2}{\bar{\tau}_k}g_{ij}
(\cdot,t_0-s\frac{\bar{\tau}_k}{2})$ and $s \in [1,+\infty)$,
converge to a manifold $(\bar{M},\bar{g}_{ij}(s),\bar{q})$ with $s
\in [1,+\infty)$, where $\bar{g}_{ij}(s)$ is also a solution to
the Ricci flow on $\bar{M}$.

Denote by $\tilde{l}_k$ the corresponding Li-Yau-Perelman distance
of $\tilde{g}_{ij}^{(k)}(s)$. It is easy to see that
$\tilde{l}_k(q,s)=l(q,\frac{\bar{\tau}_k}{2}s),$ for $s\in
[1,+\infty).$ From (6.2.5), we also have \begin{equation} |\nabla
\tilde{l}_k|^2_{\tilde{g}_{ij}^{(k)}}+\tilde{R}^{(k)} \leq
6\tilde{l}_k,
%\eqno(6.2.11)
\end{equation} where $\tilde{R}^{(k)}$ is the scalar curvature of the metric
$\tilde{g}_{ij}^{(k)}$. Claim (6.2.6) says that $\tilde{l}_k$ are
uniformly bounded on compact subsets of $M \times [1,+\infty)$
(with the corresponding origins $q(\frac{\bar{\tau}_k}{2})$). Thus
the above gradient estimate (6.2.11) implies that the functions
$\tilde{l}_k$ tend (up to a subsequence) to a function $\bar{l}$
which is a locally Lipschitz function on $\bar{M}$.

We know from (6.2.2)-(6.2.4) that the Li-Yau-Perelman distance
$\tilde{l}_k$ satisfies the following inequalities: \begin{equation}
(\tilde{l}_k)_s-\Delta
\tilde{l}_k+|\nabla\tilde{l}_k|^2-\tilde{R}^{(k)}
+\frac{n}{2s}\geq 0, %\eqno(6.2.12)
\end{equation} \begin{equation} 2\Delta\tilde{l}_k-|\nabla\tilde{l}_k|^2
+\tilde{R}^{(k)}+\frac{\tilde{l}_k-n}{s}\leq 0. %\eqno(6.2.13)
\end{equation} We next show that the limit $\bar{l}$ also satisfies the above
two inequalities in the sense of distributions. Indeed the above
two inequalities can be rewritten as \begin{equation}
\left(\frac{\partial}{\partial s}-\triangle+\widetilde{R}^{(k)}\right)
\left((4\pi s)^{-\frac{n}{2}}\exp(-\tilde{l}_k)\right)\leq 0,
%\eqno (6.2.14)
\end{equation} \begin{equation} -(4\triangle
-\widetilde{R}^{(k)})e^{-\frac{\tilde{l}_k}{2}}
+\frac{\tilde{l}_{k}-n}{s}e^{-\frac{\tilde{l}_k}{2}}\leq 0,
%\eqno(6.2.15)
\end{equation} in the sense of distributions. Note that the estimate (6.2.11)
implies that $\tilde{l}_k\rightarrow \bar{l}$ in the $C_{\rm
loc}^{0,\alpha}$ norm for any $0<\alpha<1.$ Thus the inequalities
(6.2.14) and (6.2.15) imply that the limit $\overline{l}$
satisfies \begin{equation} \left(\frac{\partial}{\partial s}-\triangle
+\overline{R}\right) \left((4\pi s)^{-\frac{n}{2}}\exp(-\overline{l})\right)\leq
0,
%\eqno(6.2.16)
\end{equation} \begin{equation} -(4\triangle-\overline{R})e^{-\frac{\overline{l}}{2}}
+\frac{\bar{l}-n}{s}e^{-\frac{\overline{l}}{2}}\leq 0,
%\eqno (6.2.17)
\end{equation} in the sense of distributions.

Denote\; by\; $\tilde{V}^{(k)}\,(s)$\; Perelman's\; reduced\;
volume\; of\; the\; scaled\; metric $\tilde{g}_{ij}^{(k)}(s)$.
Since $\tilde{l}_k(q,s)=l(q,\frac{\bar{\tau}_k}{2}s)$, we see that
$\tilde{V}^{(k)}(s)=\tilde{V}(\frac{\bar{\tau}_k}{2}s)$ where
$\tilde{V}$ is Perelman's reduced volume of the ancient
$\kappa$-solution. The monotonicity of Perelman's reduced volume
(Theorem 3.2.8) then implies that \begin{equation}
\lim_{k\rightarrow\infty}\tilde{V}^{(k)}(s)=\bar{V}, \;
\mbox{ for }\; s\in[1,2],  %\eqno(6.2.18)
\end{equation} for some nonnegative constant $\bar{V}$.

(We remark that by the Jacobian comparison theorem (Theorem
3.2.7), (3.2.18) and (3.2.19), the integrand of
$\tilde{V}^{(k)}(s)$ is bounded by
$$
(4\pi s)^{-\frac{n}{2}}\exp(-\tilde{l}_k(X,s))
\tilde{\mathcal{J}}^{(k)}(s) \leq (4\pi)^{-\frac{n}{2}}
\exp(-|X|^2)
$$
on $T_pM$, where $\tilde{\mathcal{J}}^{(k)}(s)$ is the
$\mathcal{L}$-Jacobian of the $\mathcal{L}$-exponential map of the
metric $\tilde{g}_{ij}^{(k)}(s)$ at $T_pM$. Thus we can apply the
dominant convergence theorem to get the convergence in (6.2.18).
But we are not sure whether the limiting $\bar{V}$ is exactly
Perelman's reduced volume of the limiting manifold
($\bar{M},\bar{g}_{ij}(s)$), because the points
$q(\frac{\bar{\tau}_k}{2})$ may diverge to infinity. Nevertheless,
we can ensure that $\bar{V}$ is not less than Perelman's reduced
volume of the limit.)

Note by (6.2.5) that
\begin{align}
& \tilde{V}^{(k)}(2)-\tilde{V}^{(k)}(1)\\[1mm]
& = \int_{1}^2\frac{d}{ds}(\tilde{V}^{(k)}(s))ds \nonumber\\[1mm]
& = \int_{1}^2ds\int_M\left(\frac{\partial}{\partial s}
-\Delta+\tilde{R}^{(k)}\right) \left((4\pi s)^{-\frac{n}{2}}
\exp(-\tilde{l}_k)\right)dV_{\tilde{g}_{ij}^{(k)}(s)}. \nonumber
\end{align}
Thus we deduce that in the sense of distributions, \begin{equation}
\left(\frac{\partial}{\partial s}-\Delta+\bar{R}\right) \left((4\pi
s)^{-\frac{n}{2}}\exp(-\bar{l})\right)= 0,
%\eqno (6.2.20)
\end{equation} and
$$
(4\Delta-\bar{R})e^{-\frac{\bar{l}}{2}}
=\frac{\bar{l}-n}{s}e^{-\frac{\bar{l}}{2}}
$$
or equivalently, \begin{equation} 2\Delta \bar{l} - |\nabla \bar{l}|^2 +
\bar{R}
+ \frac{\bar{l} -n}{s} = 0, %\eqno (6.2.21)
\end{equation} on $\bar{M} \times [1,2]$. Thus by applying standard parabolic
equation theory to (6.2.20) we find that $\bar{l}$ is actually
smooth. Here we used (6.2.2)-(6.2.4) to show that the equality in
(6.2.16) implies the equality in (6.2.17).

Set
$$
v=[s(2\Delta \bar{l}-|\nabla \bar{l}|^2+\bar{R})+\bar{l}-n]\cdot
(4\pi s)^{-\frac{n}{2}}e^{-\bar{l}}.
$$
Then by applying Lemma 2.6.1, we have \begin{equation}
\left(\frac{\partial}{\partial s}-\Delta+\bar{R}\right)v
=-2s|\bar{R}_{ij}+\nabla_i\nabla_j\bar{l}
-\frac{1}{2s}\bar{g}_{ij}|^2\cdot
(4\pi s)^{-\frac{n}{2}}e^{-\bar{l}}. %\eqno(6.2.22)
\end{equation} We see from (6.2.21) that the LHS of the equation (6.2.22) is
identically zero. Thus the limit metric $\bar{g}_{ij}$ satisfies
\begin{equation}
\bar{R}_{ij}+\nabla_i\nabla_j\bar{l}-\frac{1}{2s}\bar{g}_{ij}=0,
%\eqno (6.2.23)
\end{equation} so we have shown the limit is a gradient shrinking soliton.

To show that the limiting gradient shrinking soliton is nonflat,
we first show that the constant function $\bar{V}(s)$ is strictly
less than 1. Consider Perelman's reduced volume $\tilde{V}(\tau)$
of the ancient $\kappa$-solution. By using Perelman's Jacobian
comparison theorem (Theorem 3.2.7), (3.2.18) and (3.2.19) as
before, we have
\begin{align*}
\tilde{V}(\tau)&  = \int (4\pi
\tau)^{-\frac{n}{2}}e^{-l(X,\tau)}{\mathcal{J}}(\tau)dX \\
&  \leq \int_{T_pM}(4\pi)^{-\frac{n}{2}}e^{-|X|^2}dX \\
&  = 1.
\end{align*}

Recall that we have assumed the nonflat ancient $\kappa$-solution
is not a gradient shrinking soliton. Thus for $\tau>0,$ we must
have $\tilde{V}(\tau)<1$. Since the limiting function $\bar{V}(s)$
is the limit of $\tilde{V}(\frac{\bar{\tau}_k}{2}s)$ with
$\bar{\tau}_k\rightarrow +\infty$, we deduce that the constant
$\bar{V}(s)$ is strictly less than $1$, for $s\in[1,2].$

We now argue by contradiction. Suppose the limiting gradient
shrinking soliton $\bar{g}_{ij}(s)$ is flat. Then by (6.2.23),
$$
\nabla_i\nabla_j\bar{l} =\frac{1}{2s}\bar{g}_{ij} \; \mbox{ and }
\;\Delta\bar{l} =\frac{n}{2s}.
$$
Putting these into the identity (6.2.21), we get
$$
|\nabla\bar{l}|^2=\frac{\bar{l}}{s}
$$
Since the function $\bar{l}$ is strictly convex, it follows that
$\sqrt{4s\bar{l}}$ is a distance function (from some point) on the
complete flat manifold $\bar{M}$. From the smoothness of the
function $\bar{l}$, we conclude that the flat manifold $\bar{M}$
must be $\mathbb{R}^n$. In this case we would have its reduced
distance to be $\bar{l}$ and its reduced volume to be $1$. Since
$\bar{V}$ is not less than the reduced volume of the limit, this
is a contradiction. Therefore the limiting gradient shrinking
soliton $\bar{g}_{ij}$ is not flat.
%$$\eqno \#$$\vskip 0.3cm
\end{proof}

To conclude this section, we use the above theorem to derive the
classification of all two-dimensional ancient $\kappa$-solutions
which was obtained earlier by Hamilton in Section 26 of
\cite{Ha95F}.

%\vskip 0.2cm\noindent {\bf Theorem 6.2.2} \emph{
%CZ-SOURCE-LINE-14322
\begin{theorem} [classification of ancient two-dimensional kappa-solutions]
  \label{thm:cz-two-dimensional-ancient-kappa-classification}
  \dcref{Cao--Zhu Theorem 6.2.2}
  \group{Asymptotic Shrinking Solitons}
  \source{cao-zhu}

The only nonflat ancient $\kappa$-solutions to Ricci flow on
two-dimensional manifolds are the round sphere $\mathbb{S}^2$ and
the round real projective plane $\mathbb{R}\mathbb{P}^2$.
\end{theorem}

%\noindent{\bf{Proof}.}\quad
\begin{proof}
  \uses{thm:cz-ancient-kappa-asymptotic-shrinker, lem:cz-unbounded-curvature-point-selection, prop:cz-pointed-limit-at-infinity-splits-line, prop:cz-compact-shrinking-surface-soliton-rigidity, thm:cz-positive-euler-surface-convergence}
Let $g_{ij}(x,t)$ be a nonflat ancient $\kappa$-solution defined
on $M\times (-\infty,T)$ (for some $T>0$), where $M$ is a
two-dimensional manifold. Note that the ancient $\kappa$-solution
satisfies the Li-Yau-Hamilton inequality (Corollary 2.5.5). In
particular by Corollary 2.5.8, the scalar curvature of the ancient
$\kappa$-solution is pointwise nondecreasing in time. Moreover by
the strong maximum principle, the ancient $\kappa$-solution has
strictly positive curvature everywhere.

By the above Theorem 6.2.1, we know that the scalings of the
ancient $\kappa$-solution along a sequence of points $q_k$ in $M$
and a sequence of times $t_k \rightarrow -\infty$ converge to a
nonflat gradient shrinking soliton $(\bar{M},\bar{g}_{ij}(x,t))$
with $-\infty < t \leq 0$.

We first show that the limiting gradient shrinking soliton
$(\bar{M},\bar{g}_{ij}(x,t))$ has uniformly bounded curvature.
Clearly, the limiting soliton has nonnegative curvature and is
$\kappa$-noncollapsed on all scales, and its scalar curvature is
still pointwise nondecreasing in time. Thus we only need to show
that the limiting soliton has bounded curvature at $t=0$. We argue
by contradiction. Suppose the curvature of the limiting soliton is
unbounded at $t=0$. Of course in this case the limiting soliton
$\bar{M}$ is noncompact. Then by applying Lemma 6.1.4, we can
choose a sequence of points $x_j, j = 1, 2, \ldots,$ divergent to
infinity such that the scalar curvature $\bar{R}$ of the limit
satisfies
$$
\bar{R}(x_j,0) \geq j\; \mbox{ and }\; \bar{R}(x,0) \leq
4\bar{R}(x_j,0)
$$
for all $j = 1, 2,\, \ldots,$ and $x\in
B_0(x_j,{j}/{\sqrt{\bar{R}(x_{j},0)}})$. And then by the
nondecreasing (in time) of the scalar curvature, we have
$$
\bar{R}(x,t) \leq 4\bar{R}(x_j,0),
$$
for all $j = 1, 2,\, \ldots$, $x\in
B_0(x_j,{j}/{\sqrt{\bar{R}(x_{j},0)}})$ and $t\leq0$. By combining
with Hamilton's compactness theorem (Theorem 4.1.5) and the
$\kappa$-noncollapsing, we know that a subsequence of the
rescaling solutions
$$
(\bar{M},\bar{R}(x_j,0)\bar{g}_{ij}(x,t/\bar{R}(x_j,0)),x_j),
\quad j=1, 2, \ldots,
$$
converges in the $C^{\infty}_{loc}$ topology to a nonflat smooth
solution of the Ricci flow. Then Proposition 6.1.2 implies that
the new (two-dimensional) limit must be flat. This arrives at a
contradiction. So we have proved that the limiting gradient
shrinking soliton has uniformly bounded curvature.

We next adapt an argument of Perelman (cf. the proof of Lemma 1.2
in \cite{P2}) to show that the limiting soliton is compact.
Suppose the limiting soliton is (complete and) noncompact. By the
strong maximum principle we know that the limiting soliton also
has strictly positive curvature everywhere. After a shift of the
time, we may assume that the limiting soliton satisfies the
following equation \begin{equation} \nabla_i\nabla_j
f+\bar{R}_{ij}+\frac{1}{2t}\bar{g}_{ij}=0,\quad
\mbox{on }\;  -\infty<t<0, %\eqno (6.2.24)
\end{equation} everywhere for some function $f$. Differentiating the equation
(6.2.24) and switching the order of differentiations, as in the
derivation of (1.1.14), we get \begin{equation}
\nabla_i\bar{R}=2\bar{R}_{ij}\nabla_jf. %\eqno (6.2.25)
\end{equation}

Fix some $t<0$, say $t=-1$, and consider a long shortest geodesic
$\gamma(s)$, $0\leq s\leq \overline{s}$. Let $x_0=\gamma(0)$ and
$X(s)=\dot{\gamma}(s)$. Let $V(0)$ be any unit vector orthogonal
to $\dot{\gamma}(0)$ and translate $V(0)$ along $\gamma(s)$ to get
a parallel vector field $V(s)$, $0\leq s\leq \overline{s}$ on
$\gamma$. Set
$$
\widehat{V}(s)=
\begin{cases}
sV(s), & \mbox{for  $0\leq s\leq 1$},\\
V(s), & \mbox{for  $1\leq s\leq \overline{s}-1$},\\
(\overline{s}-s)V(s), & \mbox{for  $\overline{s}-1\leq s\leq
\overline{s}.$}
\end{cases}
$$
It follows from the second variation formula of arclength that
$$
\int^{\overline{s}}_0(|\dot{\widehat{V}}(s)|^2
-\bar{R}(X,\widehat{V},X,\widehat{V}))ds\geq 0.
$$
Thus we clearly have
$$
\int^{\overline{s}}_0\bar{R}(X,\widehat{V},X,\widehat{V})ds\leq
{\rm const.},
$$
and then \begin{equation} \int^{\overline{s}}_0\bar{\rm Ric}(X,X)ds\leq
{\rm const.} . %\eqno (6.2.26)
\end{equation} By integrating the equation (6.2.24) we get
$$
X(f(\gamma(\overline{s})))-X(f(\gamma(0)))
+\int^{\overline{s}}_0\bar{\rm
Ric}(X,X)ds-\frac{1}{2}\overline{s}=0
$$
and then by (6.2.26), we deduce
$$
\frac{d}{ds}(f\circ\gamma(s))\geq \frac{s}{2}-{\rm const.},
$$
$$
\mbox{and } f\circ\gamma(s)\geq \frac{s^2}{4}-{\rm const.}\cdot
s-{\rm const.}
$$
for $s>0$ large enough. Thus at large distances from the fixed
point $x_0$ the function $f$ has no critical points and is proper.
It then follows from the Morse theory that any two high level sets
of $f$ are diffeomorphic via the gradient curves of $f$. Since by
(6.2.25),
\begin{align*}
\frac{d}{ds}\bar{R}(\eta(s),-1)&  = \langle\nabla \bar{R},\dot{\eta}(s)\rangle\\
&  = 2\bar{R}_{ij}\nabla_if\nabla_jf \\
&  \geq 0
\end{align*}
for any integral curve $\eta(s)$ of $\nabla f$, we conclude that
the scalar curvature $\bar{R}(x,-1)$ has a positive lower bound on
$\bar{M}$, which contradicts the Bonnet-Myers Theorem. So we have
proved that the limiting gradient shrinking soliton is compact.

By Proposition 5.1.10, the compact limiting gradient shrinking
soliton has constant curvature. This says that the scalings of the
ancient $\kappa$-solution $(M,g_{ij}(x,t))$ along a sequence of
points $q_k \in M$ and a sequence of times $t_k \rightarrow
-\infty$ converge in the $C^{\infty}$ topology to the round
$\mathbb{S}^2$ or the round $\mathbb{R}\mathbb{P}^2$. In
particular, by looking at the time derivative of the volume and
the Gauss-Bonnet theorem, we know that the ancient
$\kappa$-solution $(M,g_{ij}(x,t))$ exists on a maximal time
interval $(-\infty,T)$ with $T < +\infty$.

Consider the scaled entropy of Hamilton \cite{Ha88}
$$
E(t)=\int_MR\log[R(T-t)]dV_t.
$$
We compute
\begin{align}
\frac{d}{dt}E(t)&  = \int_M \log[R(T-t)]\Delta
RdV_t+\int_M\left[\Delta R+R^2-\frac{R}{(T-t)}\right]dV_t\\
&  = \int_M\left[-\frac{|\nabla R|^2}{R}+R^2-rR\right]dV_t \nonumber\\
&  = \int_M\left[-\frac{|\nabla R|^2}{R}+(R-r)^2\right]dV_t \nonumber
\end{align}  %\eqno(6.2.27)
where $r=\int_MRdV_t/Vol_t(M)$ and we have used
$\Vol_t(M)=(\int_MRdV_t)\cdot(T-t)$ (by the Gauss-Bonnet theorem).

We now need an inequality of Chow in \cite{ChowE}. For sake of
completeness, we present his proof as follows. For a smooth function
$f$ on the surface $M$, one can readily check
\begin{align*}
\int_M(\Delta f)^2 &=2\int_M\left|\nabla_i\nabla_j
f-\frac{1}{2}(\Delta f)g_{ij}\right|^2
+\int_MR|\nabla f|^2, \\
\int_M\frac{|\nabla R+R\nabla f|^2}{R} &=\int_M\frac{|\nabla
R|^2}{R}-2\int_M R(\Delta f)+\int_MR|\nabla f|^2,
\end{align*}
and then
\begin{align*}
&\int_M\frac{|\nabla R|^2}{R}+\int_M(\Delta f)(\Delta f-2R)\\
&=2\int_M\left|\nabla_i\nabla_jf-\frac{1}{2}(\Delta
f)g_{ij}\right|^2 +\int_M\frac{|\nabla R+R\nabla f|^2}{R}.
\end{align*}
By choosing $\Delta f=R-r$, we get
\begin{align*}
&\int_M\frac{|\nabla R|^2}{R}-\int_M(R-r)^2 \\
&=2\int_M\left|\nabla_i\nabla_jf-\frac{1}{2}(\Delta
f)g_{ij}\right|^2 +\int_M\frac{|\nabla R+R\nabla f|^2}{R}\ge 0.
\end{align*}
If the equality holds, then we have
$$
\nabla_i\nabla_jf-\frac{1}{2}(\Delta f)g_{ij}=0
$$
i.e., $\nabla f$ is conformal. By the Kazdan-Warner identity
\cite{KW}, it follows that
$$
\int_M\nabla R\cdot \nabla f=0,
$$
so
\begin{align*}
0&  = -\int_MR\Delta f\\
&  = -\int_M(R-r)^2.
\end{align*}
Hence we have proved the following inequality due to Chow
\cite{ChowE} \begin{equation} \int_M\frac{|\nabla R|^2}{R}\ge\int_M(R-r)^2,
%\eqno(6.2.28)
\end{equation} and the equality holds if and only if $R\equiv r$.

The combination (6.2.27) and (6.2.28) shows that the scaled
entropy $E(t)$ is strictly decreasing along the Ricci flow unless
we are on the round sphere $\mathbb{S}^2$ or its quotient
$\mathbb{R}\mathbb{P}^2$. Moreover the convergence result in
Theorem 5.1.11 shows that the scaled entropy $E$ has its minimum
value at the constant curvature metric (round $\mathbb{S}^2$ or
round $\mathbb{R}\mathbb{P}^2$). We had shown that the scalings of
the nonflat ancient $\kappa$-solution along a sequence of times
$t_k\rightarrow-\infty$ converge to the constant curvature metric.
Then $E(t)$ has its minimal value at $t=-\infty$, so it was
constant all along, hence the ancient $\kappa$-solution must have
constant curvature for each $t\in (-\infty,T)$. This proves the
theorem.
%$$\eqno \#$$\vskip 1cm
\end{proof}


\section{Curvature Estimates via Volume Growth}

For solutions to the Ricci flow, Perelman's no local collapsing
theorems tell us that the local curvature upper bounds imply the
local volume lower bounds. Conversely, one would expect to get
local curvature upper bounds from local volume lower bounds. If
this is the case, one will be able to establish an elliptic type
estimate for the curvatures of solutions to the Ricci flow. This
will provide the key estimate for the canonical neighborhood
structure and thick-thin decomposition of the Ricci flow on
three-manifolds. In this section we derive such curvature
estimates for nonnegatively curved solutions. In the next chapter
we will derive similar estimates for all smooth solutions, as well
as surgically modified solutions, of the Ricci flow on
three-manifolds.

Let $M$ be an $n$-dimensional complete noncompact Riemannian
manifold with nonnegative Ricci curvature. Pick an origin $O\in
M$. The well-known Bishop-Gromov volume comparison theorem tells
us the ratio $Vol(B(O,r))/r^n$ is monotone nonincreasing in
$r\in[0,+\infty)$. Thus there exists a limit
$$
\nu_M=\lim_{r\rightarrow+\infty}\frac{\Vol(B(O,r))}{r^n}.
$$
Clearly the number $\nu_M$ is invariant under dilation and is
independent of the choice of the origin. $\nu_M$ is called the
{\bf asymptotic volume ratio}\index{asymptotic volume ratio} of
the Riemannian manifold $M$.

The following result obtained by Perelman (cf. Proposition 11.4 of
\cite{P1}) shows that any ancient $\kappa$-solution must have zero
asymptotic volume ratio. This result for the Ricci flow on
K\"ahler manifolds was implicitly and independently given by Chen
and the second author in the proof of Theorem 5.1 of \cite{CZ04}.
Moreover in the K\"ahler case, as shown by Chen, Tang and the
second author (implicitly in the proof of Theorem 4.1 of
\cite{CTZ04} for complex two dimension) and by Ni (in \cite{Ni}
for all dimensions), the condition of nonnegative curvature
operator can be replaced by the weaker condition of nonnegative
bisectional curvature.

%\vskip 0.2cm \noindent {\bf Lemma 6.3.1}\  \emph{
%CZ-SOURCE-LINE-14578
\begin{lemma} [vanishing asymptotic volume ratio of ancient solutions]
  \label{lem:cz-ancient-solution-asymptotic-volume-ratio-zero}
  \dcref{Cao--Zhu Lemma 6.3.1}
  \group{Curvature Estimates from Volume Growth}
  \source{cao-zhu}

Let $M$ be an $n$-dimensional complete noncompact Riemannian
manifold. Suppose $g_{ij}(x,t)$, $x\in M$ and $t\in(-\infty,T)$
with $T>0$, is a nonflat ancient solution of the Ricci flow with
nonnegative curvature operator and bounded curvature. Then the
asymptotic volume ratio of the solution metric satisfies
$$
\nu_M(t)=\lim_{r\rightarrow+\infty}\frac{\Vol_t(B_t(O,r))}{r^n}=0
$$
for each $t\in(-\infty,T)$.
\end{lemma}

% \vskip 0.1cm \noindent {\bf Proof.}
\begin{proof}
  \uses{thm:cz-two-dimensional-ancient-kappa-classification, lem:cz-infinite-ascr-point-selection, prop:cz-pointed-limit-at-infinity-splits-line}
The proof is by induction on the dimension. When the dimension is
two, the lemma is valid by Theorem 6.2.2. For dimension $\geq 3$,
we argue by contradiction.

Suppose the lemma is valid for dimensions $\leq n-1$ and suppose
$\nu_{M}(t_0) > 0$ for some $n$-dimensional nonflat ancient
solution with nonnegative curvature operator and bounded curvature
at some time $t_0 \leq 0$. Fixing a point $x_0 \in M$, we consider
the asymptotic scalar curvature ratio
$$
A=\limsup_{d_{t_0}(x,x_0)\rightarrow+\infty}
R(x,t_0)d^2_{t_0}(x,x_0).
$$
We divide the proof into three cases.

\medskip
{\it Case} 1: $A=+\infty$.

\smallskip
By Lemma 6.1.3, there exist sequences of points $x_k \in M$
divergent to infinity, of radii $r_k \rightarrow +\infty$, and of
positive constants $\delta_k \rightarrow 0$ such that
\begin{itemize}
\item[(i)] $R(x,t_0)\leq(1+\delta_k)R(x_k,t_0)$ for all $x$ in the
ball $B_{t_0}(x_k,r_k)$ of radius $r_k$ around $x_k$, \item[(ii)]
$r^2_kR(x_k,t_0)\rightarrow+\infty$ as $k\rightarrow+\infty$,
\item[(iii)] $d_{t_0}(x_k,x_0)/r_k\rightarrow+\infty$.
\end{itemize}

By scaling the solution around the points $x_k$ with factor
$R(x_k,t_0)$, and shifting the time $t_0$ to the new time zero, we
get a sequence of rescaled solutions
$$
{g}_k(s) = R(x_k,t_0)g\left(\cdot,t_0 + \frac{s}{R(x_k,t_0)}\right)
$$
to the Ricci flow. Since the ancient solution has nonnegative
curvature operator and bounded curvature, there holds the
Li-Yau-Hamilton inequality (Corollary 2.5.5). Thus the rescaled
solutions satisfy
$$
{R}_k(x,s) \leq (1+\delta_k)
$$
for all $s\leq 0$ and $x\in
B_{{g}_k(0)}(x_k,r_k\sqrt{R(x_k,t_0)}).$ Since $\nu_{M}(t_0)
> 0$, it follows from the standard volume comparison and Theorem
4.2.2 that the injectivity radii of the rescaled solutions ${g}_k$
at the points $x_k$ and the new time zero is uniformly bounded
below by a positive number. Then by Hamilton's compactness theorem
(Theorem 4.1.5), after passing to a subsequence,
$(M,{g}_k(s),x_k)$ will converge to a solution
$(\tilde{M},\tilde{g}(s),O)$ to the Ricci flow with
$$
\tilde{R}(y,s) \leq 1, \mbox{ for all } s \leq 0 \mbox{ and } y\in
\tilde{M},
$$
and
$$
\tilde{R}(O,0) = 1.
$$
Since the metric is shrinking, by (ii) and (iii), we get
$$
R(x_k,t_0)d^2_{g(\cdot,t_0 + \frac{s}{R(x_k,t_0)})}(x_0,x_k) \geq
R(x_k,t_0)d^2_{g(\cdot,t_0)}(x_0,x_k)
$$
which tends to $+\infty$, as $k\rightarrow +\infty$, for all
$s\leq 0$. Thus by Proposition 6.1.2, for each $s\leq 0$,
$(\tilde{M},\tilde{g}(s))$ splits off a line. We now consider the
lifting of the solution $(\tilde{M},\tilde{g}(s)), s \leq 0,$ to
its universal cover and denote it by
$(\tilde{\tilde{M}},\tilde{\tilde{g}}(s)), s \leq 0.$ Clearly we
still have
$$
\nu_{\tilde{M}}(0) > 0 \mbox{ and }\nu_{\tilde{\tilde{M}}}(0) > 0.
$$
By applying Hamilton's strong maximum principle and the de Rham
decomposition theorem, the universal cover $\tilde{\tilde{M}}$
splits isometrically as $X \times \mathbb{R}$ for some
$(n-1)$-dimensional nonflat (complete) ancient solution $X$ with
nonnegative curvature operator and bounded curvature. These imply
that $\nu_X(0) > 0$, which contradicts the induction hypothesis.

\medskip
{\it Case} 2: $0 < A < +\infty$.

\smallskip
Take a sequence of points $x_k$ divergent to infinity such that
$$
R(x_k,t_0)d^2_{t_0}(x_k,x_0) \rightarrow A,  \mbox{ as } k
\rightarrow +\infty.
$$
Consider the rescaled solutions $(M,g_k(s))$ (around the fixed
point $x_0$), where
$$
g_k(s) = R(x_k,t_0)g\left(\cdot,t_0 + \frac{s}{R(x_k,t_0)}\right), s\in
(-\infty,0].
$$
Then there is a constant $C > 0$ such that \begin{equation}
   \left\{
   \begin{array}{lll}
    R_k(x,0) \leq {C}/{d^2_k(x,x_0,0)},
         \\[2mm]
    R_k(x_k,0) = 1,
         \\[2mm]
    d_k(x_k,x_0,0) \rightarrow \sqrt{A} >0,
\end{array}
\right.%\eqno (6.3.1)
\end{equation} where $d_k(\cdot,x_0,0)$ is the distance function from the
point $x_0$ in the metric $g_k(0)$.

Since $\nu_M(t_0) >0$, it is a basic result in Alexandrov space
theory (see for example Theorem 7.6 of \cite{CC96}) that a
subsequence of $(M,g_k(s),x_0)$ converges in the Gromov-Hausdorff
sense to an $n$-dimensional metric cone
$(\tilde{M},\tilde{g}(0),x_0)$ with vertex $x_0$.

By (6.3.1), the standard volume comparison and Theorem 4.2.2, we
know that the injectivity radius of $(M,g_k(0))$ at $x_k$ is
uniformly bounded from below by a positive number $\rho_0$. After
taking a subsequence, we may assume $x_k$ converges to a point
$x_{\infty}$ in $\tilde{M}\setminus \{x_0\}$. Then by Hamilton's
compactness theorem (Theorem 4.1.5), we can take a subsequence
such that the metrics $g_k(s)$ on the metric balls
$B_0(x_k,\frac{1}{2}\rho_0) (\subset M $ with respect to the
metric $g_k(0))$ converge in the $C^{\infty}_{loc}$ topology to a
solution of the Ricci flow on a ball
$B_0(x_{\infty},\frac{1}{2}\rho_0)$. Clearly the
$C^{\infty}_{loc}$ limit has nonnegative curvature operator and it
is a piece of the metric cone at the time $s=0$. By (6.3.1), we
have \begin{equation}
\tilde{R}(x_{\infty},0) = 1. %\eqno (6.3.2)
\end{equation}

Let $x$ be any point in the limiting ball
$B_0(x_{\infty},\frac{1}{2}\rho_0)$ and $e_1$ be any radial
direction at $x$. Clearly $\tilde{Ric}(e_1,e_1) = 0$. Recall that
the evolution equation of the Ricci tensor in frame coordinates is
$$
\frac{\partial}{\partial t}{\tilde{R}_{ab}} =
\tilde{\triangle}{\tilde{R}_{ab}} +
2{\tilde{R}_{acbd}}{\tilde{R}_{cd}}.
$$
Since the curvature operator is nonnegative, by applying
Hamilton's strong maximum principle (Theorem 2.2.1) to the above
equation, we deduce that the null space of $\tilde{Ric}$ is
invariant under parallel translation. In particular, all radial
directions split off locally and isometrically. While by (6.3.2),
the piece of the metric cone is nonflat. This gives a
contradiction.

\medskip
{\it Case} 3: $A = 0$.

\smallskip
The gap theorem as was initiated by Mok-Siu-Yau \cite{MoSY} and
established by Greene-Wu \cite{GW82, GW},
Eschenberg-Shrader-Strake \cite{ESS}, and Drees \cite{Dr} shows
that a complete noncompact $n$-dimensional (except $n=4$ or $8$)
Riemannian manifold with nonnegative sectional curvature and the
asymptotic scalar curvature ratio $A=0$ must be flat. So the
present case is ruled out except in dimension $n=4$ or $8$. Since
in our situation the asymptotic volume ratio is positive and the
manifold is the solution of the Ricci flow, we can give an
alternative proof for all dimensions as follows.

We claim the sectional curvature of $(M,g_{ij}(x,t_0))$ is
positive everywhere. Indeed, by Theorem 2.2.2, the image of the
curvature operator is just the restricted holonomy algebra
$\mathcal{G}$ of the manifold. If the sectional curvature vanishes
for some two-plane, then the holonomy algebra $\mathcal{G}$ cannot
be $so(n)$. We observe the manifold is not Einstein since it is
noncompact, nonflat and has nonnegative curvature operator. If
$\mathcal{G}$ is irreducible, then by Berger's Theorem \cite{Ber},
$\mathcal{G} = u(\frac{n}{2})$. Thus the manifold is  K\"ahler
with bounded and nonnegative bisectional curvature and with
curvature decay faster than quadratic. Then by the gap theorem
obtained by Chen and the second author in \cite{CZ03}, this
K\"ahler manifold must be flat. This contradicts the assumption.
Hence the holonomy algebra $\mathcal{G}$ is reducible and the
universal cover of $M$ splits isometrically as $\tilde{M}_1 \times
\tilde{M}_2$ nontrivially. Clearly the universal cover of $M$ has
positive asymptotic volume ratio. So $\tilde{M}_1$ and
$\tilde{M}_2$ still have positive asymptotic volume ratio and at
least one of them is nonflat. By the induction hypothesis, this is
also impossible. Thus our claim is proved.

Now we know that the sectional curvature of $(M,g_{ij}(x,t_0))$ is
positive everywhere. Choose a sequence of points $x_k$ divergent
to infinity such that
$$
   \left\{
   \begin{array}{lll}
    R(x_k,t_0)d^2_{t_0}(x_k,x_0) = \sup \{ R(x,t_0)d^2_{t_0}(x,x_0)\
    |\
    d_{t_0}(x,x_0)\geq d_{t_0}(x_k,x_0)\},
         \\[2mm]
    d_{t_0}(x_k,x_0) \geq k,
         \\[2mm]
    R(x_k,t_0)d^2_{t_0}(x_k,x_0)=\varepsilon_k \rightarrow 0.
\end{array}
\right.
$$
Consider the rescaled metric
$$
g_k(0) = R(x_k,t_0)g(\cdot,t_0)
$$
as before. Then \begin{equation}
   \left\{
   \begin{array}{lll}
    R_k(x,0) \leq {\varepsilon_k}/{d^2_k(x,x_0,0)},\quad \mbox{ for }
    d_k(x,x_0,0) \geq \sqrt{\varepsilon_k},
         \\[2mm]
    d_{k}(x_k,x_0,0) = \sqrt{\varepsilon_k} \rightarrow 0.
\end{array}
\right. %\eqno (6.3.3)
\end{equation} As in Case 2, the rescaled marked solutions $(M,g_k(0),x_0)$
will converge in the Gromov-Hausdorff sense to a metric cone
$(\tilde{M},\tilde{g}(0),x_0)$. And by the virtue of Hamilton's
compactness theorem (Theorem 4.1.5), up to a subsequence, the
convergence is in the $C^{\infty}_{loc}$ topology in
$\tilde{M}\setminus \{x_0\}$. We next claim the metric cone
$(\tilde{M},\tilde{g}(0),x_0)$ is isometric to $\mathbb{R}^n$.

Indeed, let us write the metric cone $\tilde{M}$ as a warped
product $\mathbb{R}_{+}\times_r X^{n-1}$ for some
$(n-1)$-dimensional manifold $X^{n-1}$. By (6.3.3), the metric
cone must be flat and $X^{n-1}$ is isometric to a quotient of the
round sphere $\mathbb{S}^{n-1}$ by fixed point free isometries in
the standard metric. To show $\tilde{M}$ is isometric to
$\mathbb{R}^n$, we only need to verify that $X^{n-1}$ is simply
connected.

Let $\varphi$ be the Busemann function of $(M,g_{ij}(\cdot,t_0))$
with respect to the point $x_0$. Since $(M,g_{ij}(\cdot,t_0))$ has
nonnegative sectional curvature, it is easy to see that for any
small $\varepsilon >0$, there is a $r_0 >0$ such that
$$(1-\varepsilon)d_{t_0}(x,x_0) \leq \varphi(x) \leq
d_{t_0}(x,x_0)$$ for all $x \in M\setminus B_{t_0}(x_0,r_0)$. The
strict positivity of the sectional curvature of the manifold
$(M,g_{ij}(\cdot,t_0))$ implies that the square of the Busemann
function is strictly convex (and exhausting). Thus every level set
$\varphi^{-1}(a)$, with $a > \inf \{\varphi(x)\ |\ x \in M \}$, of
the Busemann function $\varphi$ is diffeomorphic to the
$(n-1)$-sphere $\mathbb{S}^{n-1}$. In particular,
$\varphi^{-1}([a,\frac{3}{2}a])$ is simply connected for $a > \inf
\{\varphi(x)\ |\ x \in M \}$ since $n \geq 3$.

Consider an annulus portion $[1,2]\times X^{n-1}$ of the metric
cone $\tilde{M} = \mathbb{R}_+\times_r X^{n-1}$. It is the limit
of $(M_k,g_k(0))$, where
$$
M_k = \left\{ x \in M \ \Big|\  \frac{1}{\sqrt{R(x_k,t_0)}} \leq
d_{t_0}(x,x_0) \leq \frac{2}{\sqrt{R(x_k,t_0)}} \right\}.
$$
It is clear that
\begin{multline*}
\varphi^{-1}\left(\left[\frac{1}{\sqrt{R(x_k,t_0)}},
\frac{2(1-\varepsilon)}{\sqrt{R(x_k,t_0)}}\right]\right)
\subset M_k \\
\subset
\varphi^{-1}\left(\left[\frac{1-\varepsilon}{\sqrt{R(x_k,t_0)}},
\frac{2}{\sqrt{R(x_k,t_0)}}\right]\right)\!\!\!
\end{multline*}
for $k$ large enough. Thus any closed loop in $\{\frac{3}{2}\}
\times X^{n-1}$ can be shrunk to a point by a homotopy in
$[1,2]\times X^{n-1}$. This shows that $X^{n-1}$ is simply
connected. Hence we have proven that the metric cone
$(\tilde{M},\tilde{g}(0),x_0)$ is isometric to $\mathbb{R}^n$.
Consequently,
$$
\lim_{k\rightarrow+\infty}\!\frac{\Vol_{g(t_0)}\!\left(\!\!B_{g(t_0)}
\left(\!x_0,\frac{r}{\sqrt{R(x_k,t_0)}}\right)\setminus
B_{g(t_0)}\!\left(\!x_0,\frac{\sigma r}{\sqrt{R(x_k,t_0)}}\right)\right)}{
\left(\frac{r}{\sqrt{R(x_k,t_0)}}\right)^n} = \alpha_n (1\!-\!\sigma^n)
$$
for any $r>0$ and $0<\sigma<1$, where $\alpha_n$ is the volume of
the unit ball in the Euclidean space $\mathbb{R}^n$. Finally, by
combining with the monotonicity of the Bishop-Gromov volume
comparison, we conclude that the manifold $(M,g_{ij}(\cdot,t_0))$
is flat and isometric to $\mathbb{R}^n$. This contradicts the
assumption.

Therefore, we have proved the lemma.
\end{proof}

Finally we would like to include an alternative simpler argument,
inspired by Ni \cite{Ni}, for the above Case 2 and Case 3 to avoid
the use of the gap theorem, holonomy groups, and asymptotic cone
structure.

\medskip
{\bf \em Alternative Proof for Case {\bf 2} and Case {\bf 3}.}

Let us consider the situation of $0 \leq A < +\infty$ in the above
proof. Observe that $\nu_M(t)$ is nonincreasing in time $t$ by
using Lemma 3.4.1(ii) and the fact that the metric is shrinking in
$t$. Suppose $\nu_M(t_0)>0$, then the solution $g_{ij}(\cdot,t)$
is $\kappa$-noncollapsed for $t\leq t_0$ for some uniform
$\kappa>0$. By combining with Theorem 6.2.1, there exist a
sequence of points $q_k$ and a sequence of times $t_k \rightarrow
-\infty$ such that the scalings of $g_{ij}(\cdot,t)$ around $q_k$
with factor $|t_k|^{-1}$ and with the times $t_k$ shifting to the
new time zero converge to a nonflat gradient shrinking soliton
$\bar{M}$ in the $C_{loc}^\infty$ topology. This gradient soliton
also has maximal volume growth (i.e. $\nu_{\bar{M}}(t)>0$) and
satisfies the Li-Yau-Hamilton estimate (Corollary 2.5.5). If the
curvature of the shrinking soliton $\bar{M}$ at the time $-1$ is
bounded, then we see from the proof of Theorem 6.2.2 that by using
the equations (6.2.24)-(6.2.26), the scalar curvature has a
positive lower bound everywhere on $\bar{M}$ at the time $-1$. In
particular, this implies the asymptotic scalar curvature ratio
$A=\infty$ for the soliton at the time $-1$, which reduces to Case
1 and arrives at a contradiction by the dimension reduction
argument. On the other hand, if the scalar curvature is unbounded,
then by Lemma 6.1.4, the Li-Yau-Hamilton estimate (Corollary
2.5.5) and Lemma 6.1.2, we can do the same dimension reduction as
in Case 1 to arrive at a contradiction also.
%$$\eqno \#$$\vskip 0.3cm
\endproof

The following lemma is a local and space-time version of Lemma
6.1.4 on picking local (almost) maximum curvature points. We
formulate it from Perelman's arguments in section 10 of \cite{P1}.

%\vskip 0.2cm \noindent {\bf Lemma 6.3.2} \emph{
%CZ-SOURCE-LINE-14916
\begin{lemma}[parabolic point-picking curvature control]
  \label{lem:cz-parabolic-point-picking-curvature-control}
  \dcref{Cao--Zhu Lemma 6.3.2}
  \group{Curvature Estimates from Volume Growth}
  \source{cao-zhu}

For any positive constants $B, C$ with $B>4$ and $C>1000$, there
exists $1\leq A<\min\{\frac{1}{4}B,\frac{1}{1000}C\}$ which tends
to infinity as $B$ and $C$ tend to infinity and satisfies the
following property. Suppose we have a (not necessarily complete)
solution $g_{ij}(t)$ to the Ricci flow, defined on $M\times
[-t_0,0]$, so that at each time $t\in [-t_0,0]$ the metric ball
$B_t(x_0,1)$ is compactly contained in $M$. Suppose there exists a
point $(x',t') \in M\times (-t_0,0]$ such that
$$
d_{t'}(x',x_0)\leq \frac{1}{4} \; \mbox{ and }\;
|Rm(x',t')|>C+B(t'+t_0)^{-1}.
$$
Then we can find a point $(\bar{x},\bar{t}) \in M\times (-t_0,0]$
such that
$$
d_{\bar{t}}(\bar{x},x_0)<\frac{1}{3}\; \mbox{ with }\; Q
=|Rm(\bar{x},\bar{t})|>C+ B(\bar{t}+t_0)^{-1},
$$
and
$$
|Rm(x,t)|\leq 4Q
$$
for all $(-t_0<) \ \bar{t}-AQ^{-1}\leq t\leq \bar{t}$ and
$d_t(x,\bar{x})\leq \frac{1}{10}A^{\frac{1}{2}}Q^{-\frac{1}{2}}.$
\end{lemma}

%\vskip 0.1cm \noindent{\bf Proof.}
\begin{proof}
  \uses{lem:cz-finite-parabolic-curvature-point-selection}
We first claim that there exists a point $(\bar{x},\bar{t})$ with
$-t_0<\bar{t}\leq 0$ and $d_{\bar{t}}(\bar{x},x_0)<\frac{1}{3}$
such that
$$
Q=|Rm(\bar{x},\bar{t})|>C+ B(\bar{t}+t_0)^{-1},
$$
and \begin{equation}
|Rm(x,t)|\leq 4Q %\eqno (6.3.1)? ---> 6.3.4     WRONG number
\end{equation} wherever $\bar{t}-AQ^{-1}\leq t\leq \bar{t},\ \ d_t(x,x_0)\leq
d_{\bar{t}}(\bar{x},x_0) +(AQ^{-1})^{\frac{1}{2}}.$

We will construct such $(\bar{x},\bar{t})$ as a limit of a finite
sequence of points.  Take an arbitrary $(x_1,t_1)$ such that
$$
d_{t_1}(x_1,x_0)\leq \frac{1}{4},\ \ \ -t_0<t_1\leq 0 \; \mbox{
and } \;|Rm(x_1,t_1)|>C+B(t_1+t_0)^{-1}.
$$
Such a point clearly exists by our assumption. Assume we have
already constructed $(x_k,t_k).$ If we cannot take the point
$(x_k,t_k)$ to be the desired point $(\bar{x},\bar{t})$, then
there exists a point $(x_{k+1},t_{k+1})$ such that
$$
t_k-A|Rm(x_k,t_k)|^{-1}\leq t_{k+1}\leq t_k,
$$
and
$$
d_{t_{k+1}}(x_{k+1},x_0) \leq
d_{t_{k}}(x_{k},x_0)+(A|Rm(x_k,t_k)|^{-1})^{\frac{1}{2}},
$$
but
$$
|Rm(x_{k+1},t_{k+1})|>4|Rm(x_k,t_k)|.$$

\vskip 0.1cm\noindent It then follows that
\begin{align*}
d_{t_{k+1}}(x_{k+1},x_0)& \leq
d_{t_{1}}(x_{1},x_0)+A^{\frac{1}{2}}
\left(\sum\limits_{i=1}^k|Rm(x_i,t_i)|^{-\frac{1}{2}}\right)\\
&  \leq \frac{1}{4}+A^{\frac{1}{2}}
\left(\sum\limits_{i=1}^k2^{-(i-1)}|Rm(x_1,t_1)|^{-\frac{1}{2}}\right)\\
&  \leq \frac{1}{4}+2(AC^{-1})^{\frac{1}{2}} \\
&  < \frac{1}{3},
\end{align*}
\begin{displaymath}
\begin{split}
t_{k+1}-(-t_0)&  =\sum_{i=1}^k(t_{i+1}-t_i)+(t_1-(-t_0))\\[1mm]
&  \geq-\sum_{i=1}^kA|Rm(x_i,t_i)|^{-1}+(t_1-(-t_0))\\[1mm]
&  \geq-A\sum_{i=1}^k4^{-(i-1)}|Rm(x_1,t_1)|^{-1}+(t_1-(-t_0))\\[1mm]
&  \geq-\frac{2A}{B}(t_1+t_0)+(t_1+t_0)\\[1mm]
&  \geq\frac{1}{2}(t_1+t_0),
\end{split}
\end{displaymath}
and
\begin{align*}
|Rm(x_{k+1},t_{k+1})|&  > 4^k|Rm(x_1,t_1)| \\
&  \geq 4^kC\rightarrow +\infty\; \mbox{ as } \;
k\rightarrow+\infty.
\end{align*}
Since the solution is smooth, the sequence $\{(x_k, t_k)\}$ is
finite and its last element fits. Thus we have proved  assertion
(6.3.4).

{}From the above construction we also see that the chosen point
$(\bar{x},\bar{t})$ satisfies
$$
d_{\bar{t}}(\bar{x},x_0) < \frac{1}{3}
$$
and
$$
Q=|Rm(\bar{x},\bar{t})|>C+B(\bar{t}+t_0)^{-1}.
$$
Clearly, up to some adjustment of the constant $A$, we only need
to show that
$$
|Rm(x,t)|\leq 4Q\leqno(6.3.4)' %\eqno(6.3.2)--->(6.3.4)' WRONG number
$$
wherever $\bar{t}-\frac{1}{200n}A^{\frac{1}{2}}Q^{-1}\leq t\leq
\bar{t} \ \ \mbox{and} \ \ d_t(x,\bar{x})\leq
\frac{1}{10}A^{\frac{1}{2}}Q^{-\frac{1}{2}}$.

For any point $(x,\bar{t})$ with $d_{\bar{t}}(x,\bar{x})\leq
\frac{1}{10}A^{\frac{1}{2}}Q^{-\frac{1}{2}}$, we have
\begin{align*}
d_{\bar{t}}(x,x_0)&  \leq
d_{\bar{t}}(\bar{x},x_0)+d_{\bar{t}}(x,\bar{x})\\
&  \leq d_{\bar{t}}(\bar{x},x_0)+(AQ^{-1})^{\frac{1}{2}}
\end{align*}
and then by (6.3.4)
$$
|Rm(x,\bar{t})|\leq 4Q.
$$
Thus by continuity, there is a minimal
$\bar{t}'\in[\bar{t}-\frac{1}{200n}A^{\frac{1}{2}}Q^{-1},\bar{t}]$
such that \begin{equation} \sup\left\{|Rm(x,t)|\  | \  \bar{t}'\leq t\leq
\bar{t},\quad d_t(x,\bar{x})\leq
\frac{1}{10}A^{\frac{1}{2}}Q^{-\frac{1}{2}}\right\}\leq 5Q.
%\eqno(6.3.3) ---> (6.3.5)  WRONG number
\end{equation}

For any point $(x,t)$ with $\bar{t}'\leq t\leq \bar{t}$ and
$d_t(x,\bar{x})\leq \frac{1}{10}(AQ^{-1})^{\frac{1}{2}}$, we
divide the discussion into two cases.

\medskip
{\it Case} (1): $d_t(\bar{x},x_0)\leq
\frac{3}{10}(AQ^{-1})^{\frac{1}{2}}$.

\smallskip
{}From assertion (6.3.4) we see that
$$
\sup\{|Rm(x,t)|\ \ | \ \ \bar{t}'\leq t\leq \bar{t},\ \
d_t(x,x_0)\leq (AQ^{-1})^{\frac{1}{2}}\}\leq 4Q. \leqno (6.3.5)'
%\eqno (6.3.4)? ---> (6.3.5)'  WRONG number
$$
Since $d_t(\bar{x},x_0)\leq \frac{3}{10}(AQ^{-1})^{\frac{1}{2}}$,
we have
\begin{align*}
d_t({x},x_0) &\leq d_t(x,\bar{x})+d_t(\bar{x},x_0)\\
&  \leq \frac{1}{10}(AQ^{-1})^{\frac{1}{2}}
+\frac{3}{10}(AQ^{-1})^{\frac{1}{2}}\\
&  \leq (AQ^{-1})^{\frac{1}{2}}
\end{align*}
which implies the estimate $|Rm(x,t)|\leq 4Q $ from (6.3.5)$'$.

\medskip
{\it Case} (2): $d_t(\bar{x},x_0)>
\frac{3}{10}(AQ^{-1})^{\frac{1}{2}}$.

\smallskip
{}From the curvature bounds in (6.3.5) and (6.3.5)$'$, we can
apply Lemma 3.4.1 (ii) with $r_0=\frac{1}{10}Q^{-\frac{1}{2}}$ to
get
$$
\frac{d}{dt}(d_t(\bar{x},x_0))\geq -40(n-1)Q^{\frac{1}{2}}
$$
and then
\begin{align*}
d_t(\bar{x},x_0)&  \leq d_{\hat{t}}(\bar{x},x_0)
+40n(Q^{\frac{1}{2}})\left(\frac{1}{200n}A^{\frac{1}{2}}Q^{-1}\right)\\
& \leq d_{\hat{t}}(\bar{x},x_0)+\frac{1}{5}(AQ^{-1})^{\frac{1}{2}}
\end{align*}
where $\hat{t} \in (t,\bar{t}]$ satisfies the property that
$d_s(\bar{x},x_0)\geq \frac{3}{10}(AQ^{-1})^{\frac{1}{2}}$
whenever $s \in [t,\hat{t}]$. So we have either
\begin{align*}
d_t({x},x_0)&  \leq d_t(x,\bar{x})+d_t(\bar{x},x_0)\\
&  \leq \frac{1}{10}(AQ^{-1})^{\frac{1}{2}}
+\frac{3}{10}(AQ^{-1})^{\frac{1}{2}}+\frac{1}{5}(AQ^{-1})^{\frac{1}{2}}\\
&  \leq (AQ^{-1})^{\frac{1}{2}},
\end{align*}
or
\begin{align*}
d_t(x,x_0)&  \leq d_t(x,\bar{x})+d_t(\bar{x},x_0)\\
&  \leq \frac{1}{10}(AQ^{-1})^{\frac{1}{2}}
+d_{\bar{t}}(\bar{x},x_0)+\frac{1}{5}(AQ^{-1})^{\frac{1}{2}} \\
&  \leq d_{\bar{t}}(\bar{x},x_0)+(AQ^{-1})^{\frac{1}{2}}.
\end{align*}
It then follows from (6.3.4) that $|Rm(x,t)|\leq 4Q $.

Hence we have proved
$$
|Rm(x,t)|\leq 4Q
$$
for any point $(x,t)$ with $\bar{t}'\leq t\leq \bar{t}$ and
$d_t(x,\bar{x})\leq \frac{1}{10}(AQ^{-1})^{\frac{1}{2}}$. By
combining with the choice of $\bar{t}'$ in (6.3.5), we must have
$\bar{t}'= \bar{t}-\frac{1}{200n}A^{\frac{1}{2}}Q^{-1}$. This
proves assertion (6.3.4)$'$.

Therefore we have completed the proof of the lemma.
%$$\eqno \#$$\vskip 0.3cm
\end{proof}

We now use the volume lower bound assumption to establish the
crucial curvature upper bound estimate of Perelman \cite{P1} for
the Ricci flow. For the Ricci flow on K\"ahler manifolds, a global
version of this estimate (i.e., curvature decaying linear in time
and quadratic in space) was independently obtained in \cite{CTZ04}
and \cite{CZ04}. Note that the volume estimate conclusion in the
following Theorem 6.3.3 (ii) was not stated in Corollary 11.6 (b)
of Perelman \cite{P1}. The estimate will be used later in the
proof of Theorem 7.2.2 and Theorem 7.5.2.

%\vskip 0.2cm \noindent{\bf Theorem 6.3.3}
%CZ-SOURCE-LINE-15129
\begin{theorem}[curvature estimate from local volume growth]
  \label{thm:cz-local-volume-growth-curvature-estimate}
  \dcref{Cao--Zhu Theorem 6.3.3}
  \group{Curvature Estimates from Volume Growth}
  \source{cao-zhu}

For every $w>0$ there exist $B=B(w)<+\infty,\ \ C=C(w)<+\infty, \
\ \tau_0=\tau_0(w)>0,$ and $\xi=\xi(w)>0$ $($depending also on the
dimension$)$ with the following properties. Suppose we have a
$($not necessarily complete$)$ solution $g_{ij}(t)$ to the Ricci
flow, defined on $M\times[-t_0r_0^2,0],$ so that at each time
$t\in [-t_0r_0^2,0]$ the metric ball $B_t(x_0,r_0)$ is compactly
contained in $M.$
\begin{itemize}
\item[(i)] If at each time $t \in [-t_0r_0^2,0]$,
$$
Rm(.,t)\geq -r_0^{-2} \; \mbox{ on }\; B_t(x_0,r_0)
$$
$$
\mbox{and } \; \Vol_t(B_t(x_0,r_0))\geq wr_0^n,
$$
then we have the estimate
$$
|Rm(x,t)|\leq Cr_0^{-2}+B(t+t_0r_0^2)^{-1}
$$
whenever $-t_0r_0^2 < t\leq 0$ and $d_t(x,x_0)\leq
\frac{1}{4}r_0.$ \item[(ii)] If for some $0 < \bar{\tau} \leq
t_0$,
$$
Rm(x,t)\geq -r_0^{-2} \; \mbox{ for }\; t\in [-\bar{\tau}
r_0^2,0], x\in B_t(x_0,r_0),
$$
$$
\mbox{and } \; \Vol_0(B_0(x_0,r_0))\geq wr_0^n,
$$
then we have the estimates
$$
\Vol_t(B_t(x_0,r_0))\geq \xi r_0^n\; \mbox{ for all }\;
\max\{-\bar{\tau}r_0^2,-\tau_0r_0^2\}   \leq t \leq 0,
$$
and
$$
|Rm(x,t)|\leq Cr_0^{-2}+B(t -\max\{-\bar{\tau}r_0^2,-\tau_0r_0^2\}
)^{-1}
$$
whenever $\max\{-\bar{\tau}r_0^2,-\tau_0r_0^2\} < t\leq 0$ and
$d_t(x,x_0)\leq \frac{1}{4}r_0.$
\end{itemize}
\end{theorem}

%\vskip 0.1cm \noindent{\bf Proof.}
\begin{proof}
  \uses{lem:cz-parabolic-point-picking-curvature-control, lem:cz-ancient-solution-asymptotic-volume-ratio-zero}
By scaling we may assume $r_0=1.$

\medskip
(i) By the standard (relative) volume comparison, we know that
there exists some $w'>0$, with $w'\leq w $, depending only on $w$,
such that for each point $(x,t)$ with $-t_0\leq t\leq 0$ and
$d_t(x,x_0)\leq \frac{1}{3},$ and for each $r\leq \frac{1}{3}$,
there holds \begin{equation}
\Vol_t(B_t(x,r))\geq w'r^n. %\eqno (6.3.6)  WRONG
\end{equation}

We argue by contradiction. Suppose there are sequences $B,
C\rightarrow +\infty,$ of solutions $g_{ij}(t)$ and points
$(x',t')$ such that
$$
d_{t'}(x',x_0)\leq \frac{1}{4},\quad -t_0<t'\leq 0\; \mbox{ and }
\; |Rm(x',t')|>C+B(t'+t_0)^{-1}.
$$
Then by Lemma 6.3.2, we can find a sequence of points
$(\bar{x},\bar{t})$ such that
$$
d_{\bar{t}}(\bar{x},x_0)< \frac{1}{3},
$$
$$
Q=|Rm(\bar{x},\bar{t})|> C+B(\bar{t}+t_0)^{-1},
$$
and
$$
|Rm(x,t)|\leq 4Q
$$
wherever $(-t_0<) \ \bar{t}-AQ^{-1}\leq t\leq \bar{t},\ \
d_{t}(x,\bar{x})\leq \frac{1}{10}A^{\frac{1}{2}}Q^{-\frac{1}{2}}$,
where $A$ tends to infinity with $B,C$. Thus we may take a blow-up
limit along the points $(\bar{x},\bar{t})$ with factors $Q$ and
get a non-flat ancient solution
$(M_{\infty},g_{ij}^{(\infty)}(t))$ with nonnegative curvature
operator and with the asymptotic volume ratio
$\nu_{M_{\infty}}(t)\geq w'>0$ for each $t\in(-\infty,0] \ \ $(by
(6.3.6)). This contradicts  Lemma 6.3.1.

\medskip
(ii) Let $B(w),\ C(w)$ be good for the first part of the theorem.
By the volume assumption at $t=0$ and the standard (relative)
volume comparison, we still have the estimate
$$
\Vol_0(B_0(x,r))\geq w'r^n \leqno{(6.3.6)'}  %% WRONG
$$
for each $x \in M$ with $d_0(x,x_0)\leq \frac{1}{3}$ and $r\leq
\frac{1}{3}$. We will show that $\xi=5^{-n}w'$, $B=B(5^{-n}w')$
and $C=C(5^{-n}w')$ are good for the second part of the theorem.

By continuity and the volume assumption at $t=0$, there is a
maximal subinterval $[-\tau,0]$ of the time interval
$[-\bar{\tau},0]$ such that
$$
\Vol_t(B_t(x_0,1))\geq w \geq 5^{-n}w'\quad \mbox{for all } \;
t\in[-\tau,0].
$$
This says that the assumptions of (i) hold with $5^{-n}w'$ in
place of $w$ and with $\tau$ in place of $t_0$. Thus the
conclusion of the part (i) gives us the estimate \begin{equation}
|Rm(x,t)|\leq C+B(t+\tau)^{-1} %\eqno(6.3.7)%% WRONG
\end{equation} whenever $t\in (-\tau,0]$ and $d_t(x,x_0)\leq \frac{1}{4}$.

We need to show that one can choose a positive $\tau_0$ depending
only on $w$ and the dimension such that the maximal $\tau \geq
\min \{\bar{\tau},\tau_0\}$.

For $t\in(-\tau,0]$ and $\frac{1}{8}\leq d_t(x,x_0)\leq
\frac{1}{4}$, we use (6.3.7) and Lemma 3.4.1(ii) to get
$$
\frac{d}{dt}d_t(x,x_0)\geq
-10(n-1)(\sqrt{C}+(\sqrt{B}/\sqrt{t+\tau}))
$$
which further gives
$$
d_0(x,x_0)\geq
d_{-\tau}(x,x_0)-10(n-1)(\tau\sqrt{C}+2\sqrt{B\tau}).
$$
This means \begin{equation} B_{(-\tau)}(x_0,\frac{1}{4})\supset
B_0\left(x_0,\frac{1}{4}-10(n-1)(\tau\sqrt{C}+2\sqrt{B\tau})\right).
%\eqno(6.3.8)%% WRONG
\end{equation}

Note that the scalar curvature $R\geq -C(n)$ for some constant
$C(n)$ depending only on the dimension since $Rm\geq -1.$ We have
\begin{align*}
&\frac{d}{dt}\Vol_t\left(B_0\left(x_0,\frac{1}{4}
-10(n-1)(\tau\sqrt{C}+2\sqrt{B\tau})\right)\right)\\
&  = \int_{B_0(x_0,\frac{1}{4}-10(n-1)(\tau\sqrt{C}
+2\sqrt{B\tau}))}(-R)dV_t\\
&  \leq C(n)\Vol_t\left(B_0\left(x_0,\frac{1}{4}
-10(n-1)(\tau\sqrt{C}+2\sqrt{B\tau})\right)\right)
\end{align*}
and then
\begin{align}
&  \Vol_t\left(B_0\left(x_0,\frac{1}{4}-10(n-1)(\tau\sqrt{C}
+2\sqrt{B\tau})\right)\right)\\
& \leq e^{C(n)\tau}\Vol_{(-\tau)}\left(B_0\left(x_0,\frac{1}{4}
-10(n-1)(\tau\sqrt{C}+2\sqrt{B\tau})\right)\right).\nonumber
\end{align}  %%\eqno (6.3.9)  WRONG
Thus by (6.3.6)$'$, (6.3.8) and (6.3.9),
\begin{align*}
& \Vol_{(-\tau)}(B_{(-\tau)})(x_0,1) \\
&  \geq \Vol_{(-\tau)}(B_{(-\tau)})\left(x_0,\frac{1}{4}\right)\\
&  \geq \Vol_{(-\tau)}\left(B_0\left(x_0,\frac{1}{4}-10(n-1)(\tau\sqrt{C}
+2\sqrt{B\tau})\right)\right)\\
&  \geq e^{-C(n)\tau}\Vol_{0}\left(B_0\left(x_0,\frac{1}{4}
-10(n-1)(\tau\sqrt{C}+2\sqrt{B\tau})\right)\right)\\
& \geq e^{-C(n)\tau}w'\left(\frac{1}{4}
-10(n-1)(\tau\sqrt{C}+2\sqrt{B\tau})\right)^n.
\end{align*}
So it suffices to choose $\tau_0=\tau_0(w)$ small enough so that
$$
e^{-C(n)\tau_0}\left(\frac{1}{4}-10(n-1)(\tau_0\sqrt{C}
+2\sqrt{B\tau_0})\right)^n\geq \left(\frac{1}{5}\right)^n.
$$
Therefore we have proved the theorem.
%$$\eqno \#$$\vskip 1cm
\end{proof}


\section{Ancient \texorpdfstring{$\kappa$}{kappa}-solutions on Three-manifolds}

In this section we will determine the structures of ancient
$\kappa$-solutions on three-manifolds.

First of all, we consider a special class of ancient solutions ---
gradient shrinking Ricci solitons. Recall that a solution
$g_{ij}(t)$ to the Ricci flow is said to be a \textbf{gradient
shrinking Ricci soliton}\index{gradient shrinking Ricci soliton}
if there exists a smooth function $f$ such that \begin{equation}
\nabla_i\nabla_j f+R_{ij}+\frac{1}{2t}g_{ij}=0\; \mbox{ for }\;
-\infty<t<0. %\eqno (6.4.1)
\end{equation} A gradient shrinking Ricci soliton moves by the one parameter
group of diffeomorphisms generated by $\nabla f$ and shrinks by a
factor at the same time.

The following result of Perelman \cite{P2} gives a complete
classification for all three-dimensional complete
$\kappa$-noncollapsed gradient shrinking solitons with bounded and
nonnegative sectional curvature.

%\vskip 0.2cm\noindent{\bf Lemma 6.4.1}
%CZ-SOURCE-LINE-15320
\begin{lemma}[classification of three-dimensional shrinking solitons]
  \label{lem:cz-three-dimensional-shrinking-soliton-classification}
  \dcref{Cao--Zhu Lemma 6.4.1}
  \group{Three-Dimensional Ancient Kappa-Solutions}
  \source{cao-zhu}
\index{classification of three-dimensional shrinking
solitons} Let $(M,g_{ij}(t))$ be a nonflat gradient shrinking
soliton on a three-manifold. Suppose $(M,g_{ij}(t))$ has bounded
and nonnegative sectional curvature and is $\kappa$-noncollapsed
on all scales for some $\kappa>0$. Then $(M,g_{ij}(t))$ is one of
the following:
\begin{itemize}
\item[(i)] the round three-sphere $\mathbb{S}^3$, or a metric
quotient of $\mathbb{S}^3$; \item[(ii)] the round infinite
cylinder $\mathbb{S}^2\times \mathbb{R}$, or one of its
$\mathbb{Z}_2$ quotients.
\end{itemize}
\end{lemma}

%\noindent{\textbf{Proof}.}
\begin{proof}
  \uses{thm:cz-three-manifold-positive-ricci-sphere, rem:cz-differentiable-sphere-subsequential-convergence, prop:cz-pointed-limit-at-infinity-splits-line, thm:cz-two-dimensional-ancient-kappa-classification}
We first consider the case that the sectional curvature of the
nonflat gradient shrinking soliton is not strictly positive. Let
us pull back the soliton to its universal cover. Then the
pull-back metric is again a nonflat ancient $\kappa$-solution. By
Hamilton's strong maximum principle (Theorem 2.2.1), we know that
the pull-back solution splits as the metric product of a
two-dimensional nonflat ancient $\kappa$-solution and
$\mathbb{R}$. Since the two-dimensional nonflat ancient
$\kappa$-solution is simply connected, it follows from Theorem
6.2.2 that it must be the round sphere $\mathbb{S}^2$. Thus, the
gradient shrinking soliton must be $\mathbb{S}^2\times\mathbb{R}
/\Gamma$, a metric quotient of the round cylinder.

For each $\sigma \in \Gamma$ and $(x,s) \in
\mathbb{S}^2\times\mathbb{R}$, we write $\sigma(x,s) =
(\sigma_1(x,s),\sigma_2(x,s))$ $\in \mathbb{S}^2\times\mathbb{R}$.
Since $\sigma$ sends lines to lines, and sends cross spheres to
cross spheres, we have $\sigma_2(x,s) = \sigma_2(y,s)$, for all
$x, y \in \mathbb{S}^2$. This says that $\sigma_2$ reduces to a
function of $s$ alone on $\mathbb{R}$. Moreover, for any $(x,s),
(x',s') \in \mathbb{S}^2\times\mathbb{R}$, since $\sigma$
preserves the distances between cross spheres $\mathbb{S}^2 \times
\{s\}$ and $\mathbb{S}^2 \times \{s'\}$, we have $|\sigma_2(x,s) -
\sigma_2(x',s')| = |s-s'|$. So the projection $\Gamma_2$ of
$\Gamma$ to the second factor $\mathbb{R}$ is an isometry subgroup
of $\mathbb{R}$. If the metric quotient
$\mathbb{S}^2\times\mathbb{R} /\Gamma$
 were compact, it would not be $\kappa$-noncollapsed on sufficiently
large scales as $t \rightarrow -\infty$. Thus the metric quotient
$\mathbb{S}^2\times\mathbb{R} /\Gamma$ is noncompact. It follows
that $\Gamma_2 = \{1\} \mbox{ or } \mathbb{Z}_2$. In particular,
there is a $\Gamma$-invariant cross sphere $\mathbb{S}^2$ in the
round cylinder $\mathbb{S}^2\times\mathbb{R}$. Denote it by
$\mathbb{S}^2\times \{0\}$. Then $\Gamma$ acts on the round
two-sphere $\mathbb{S}^2\times \{0\}$ isometrically without fixed
points. This implies $\Gamma$ is either $\{1\}$ or $\mathbb{Z}_2$.
Hence we conclude that the gradient shrinking soliton is either
the round cylinder $\mathbb{S}^2\times\mathbb{R}$, or
$\mathbb{R}\mathbb{P}^2 \times \mathbb{R}$, or the twisted product
$\mathbb{S}^2 \tilde{\times} \mathbb{R}$ where $\mathbb{Z}_2$
flips both $\mathbb{S}^2$ and $\mathbb{R}$.

We next consider the case that the gradient shrinking soliton is
compact and has strictly positive sectional curvature everywhere.
By the proof of Theorem 5.2.1 (see also Remark 5.2.8) we see that
the compact gradient shrinking soliton is getting round and tends
to a space form (with positive constant curvature) as the time
approaches the maximal time $t=0$. Since the shape of a gradient
shrinking Ricci soliton does not change up to reparametrizations
and homothetical scalings, the gradient shrinking soliton has to
be the round three-sphere $\mathbb{S}^3$ or a metric quotient of
$\mathbb{S}^3$.

Finally we want to exclude the case that the gradient shrinking
soliton is noncompact and has strictly positive sectional
curvature everywhere. The following argument follows the proof of
Lemma 1.2 of Perelman \cite{P2}.

Suppose there is a (complete three-dimensional) noncompact
$\kappa$-non\-collapsed gradient shrinking soliton $g_{ij}(t)$,
$-\infty<t<0$, with bounded and positive sectional curvature at
each $t\in (-\infty,0)$ and satisfying the equation (6.4.1). Then
as in (6.2.25), we have \begin{equation}
\nabla_iR=2R_{ij}\nabla_jf. %\eqno(6.4.2)
\end{equation}

Fix some $t<0$, say $t=-1$, and consider a long shortest geodesic
$\gamma(s)$, $0\le s\le\bar{s}.$ Let $x_0=\gamma(0)$ and
$X(s)=\dot{\gamma}(s).$ Let $U(0)$ be any unit vector orthogonal
to $\dot{\gamma}(0)$ and translate $U(0)$ along $\gamma(s)$ to get
a parallel vector field $U(s)$, $0\le s\le\bar{s}$, on $\gamma$.
Set
$$
\widetilde{U}(s)=\left\{\arraycolsep=1.5pt\begin{array}{lll}
&  sU(s),\quad &  \mbox{for}\quad 0\le s\le 1,\\[2mm]
&  U(s),\quad &  \mbox{for}\quad 1\le s\le\bar{s}-1\\[2mm]
&  (\bar{s}-s)U(s),\quad &  \mbox{for}\quad \bar{s}-1\le s\le
\bar{s}.\end{array}\right.
$$
It follows from the second variation formula of arclength that
$$
\int_0^{\bar{s}}(|\dot{\widetilde{U}}(s)|^2
-R(X,\widetilde{U},X,\widetilde{U}))ds\ge 0.
$$
Since the curvature of the metric $g_{ij}(-1)$ is bounded, we
clearly have
$$
\int_0^{\bar{s}}R(X,U,X,U)ds\le {\rm const.}
$$
and then \begin{equation}
\int_0^{\bar{s}}\Ric(X,X)ds\le {\rm const.}. %\eqno(6.4.3)
\end{equation} Moreover, since the curvature of the metric $g_{ij}(-1)$ is
positive, it follows from the Cauchy-Schwarz inequality that for
any unit vector field $Y$ along $\gamma$ and orthogonal to
$X(=\dot{\gamma}(s)),$ we have
\begin{align*}
\int_0^{\bar{s}}|\Ric(X,Y)|^2ds
&  \le\int_0^{\bar{s}}\Ric(X,X)\Ric(Y,Y)ds\\
&  \le {\rm const.}\,\cdot\int_0^{\bar{s}}\Ric(X,X)ds\\
&  \le {\rm const.}
\end{align*}
and then \begin{equation} \int_0^{\bar{s}}|\Ric(X,Y)|ds
\le {\rm const.}\,\cdot(\sqrt{\bar{s}}+1). %\eqno(6.4.4)
\end{equation} {}From (6.4.1) we know
$$
\nabla_X\nabla_Xf+\Ric(X,X)-\frac{1}{2}=0
$$
and by integrating this equation we get
$$
X(f(\gamma(\bar{s})))-X(f(\gamma(0)))
+\int_0^{\bar{s}}\Ric(X,X)ds-\frac{1}{2}\bar{s}=0.
$$
Thus by (6.4.3) we deduce \begin{equation} \frac{\bar{s}}{2}-{\rm const.}\,\le
\langle X,\nabla f(\gamma(\bar{s}))\rangle
\le\frac{\bar{s}}{2}+{\rm const.}.  %\eqno(6.4.5)
\end{equation} Similarly by integrating (6.4.1) and using (6.4.4) we can
deduce \begin{equation} |\langle Y,\nabla f(\gamma(\bar{s}))\rangle|
\le {\rm const.}\,\cdot (\sqrt{\bar{s}}+1). %\eqno(6.4.6)
\end{equation} These two inequalities tell us that at large distances from
the fixed point $x_0$ the function $f$ has no critical point, and
its gradient makes a small angle with the gradient of the distance
function from $x_0$.

Now from (6.4.2) we see that at large distances from $x_0$, $R$ is
strictly increasing along the gradient curves of $f$, in
particular
$$
\bar{R}=\limsup_{d_{(-1)}(x,x_0)\rightarrow+\infty}R(x,-1)>0.
$$
Let us choose a sequence of points $(x_k,-1)$ where
$R(x_k,-1)\rightarrow\bar{R}$. By the noncollapsing assumption we
can take a limit along this sequence of points of the gradient
soliton and get an ancient $\kappa$-solution defined on
$-\infty<t<0$. By Proposition 6.1.2, we deduce that the limiting
ancient $\kappa$-solution splits off a line. Since the soliton has
positive sectional curvature, we know from Gromoll-Meyer \cite{GM}
that it is orientable. Then it follows from Theorem 6.2.2 that the
limiting solution is the shrinking round infinite cylinder with
scalar curvature $\bar{R}$ at time $t=-1$. Since the limiting
solution exists on $(-\infty,0)$, we conclude that $\bar{R}\le 1$.
Hence
$$
R(x,-1)< 1
$$
when the distance from $x$ to the fixed $x_0$ is large enough on
the gradient shrinking soliton.

Let us consider the level surface $\{f=a\}$ of $f$. The second
fundamental form of the level surface is given by
\begin{align*}
h_{ij}&  = \left\langle\nabla_i\left(\frac{\nabla f}{|\nabla
f|}\right),e_j\right\rangle\\
&  = \nabla_i\nabla_j f/|\nabla f|,\qquad i,j=1,2,
\end{align*}
where $\{e_1,e_2\}$ is an orthonormal basis of the level surface.
By (6.4.1), we have
$$
\nabla_{e_i}\nabla_{e_i}f =\frac{1}{2}-\Ric(e_i,e_i)
\ge\frac{1}{2}-\frac{R}{2}>0,\qquad i=1,2,
$$
since for a three-manifold the positivity of sectional curvature
is equivalent to $R\ge 2\Ric$. It then follows from the first
variation formula that
\begin{align}
\frac{d}{da}{\rm Area}\,\{f=a\}
& = \int_{\{f=a\}}{\rm div}\,\left(\frac{\nabla f}{|\nabla f|}\right)\\
&  \ge \int_{\{f=a\}}\frac{1}{|\nabla f|}(1-R)\nonumber\\
&  > \int_{\{f=a\}}\frac{1}{|\nabla f|}(1-\bar{R})\nonumber\\
&  \ge 0\nonumber
\end{align}     %\eqno(6.4.7)
for $a$ large enough. We conclude that Area $\{f=a\}$ strictly
increases as $a$ increases. From (6.4.5) we see that for $s$ large
enough
$$
\left|\frac{df}{ds}-\frac{s}{2}\right|\le{\rm const.},
$$
and then
$$
\left|f-\frac{s^2}{4}\right|\le {\rm const.}\,\cdot(s+1).
$$
Thus we get from (6.4.7)
$$
\frac{d}{da}{\rm Area}\,\{f=a\}
>\frac{1-\bar{R}}{2\sqrt{a}}{\rm Area}\,\{f=a\}
$$
for $a$ large enough. This implies that
$$
\log {\rm Area}\,\{f=a\}>(1-\bar{R})\sqrt{a}-{\rm const.}
$$
for $a$ large enough. But it is clear from (6.4.7) that
Area$\,\{f=a\}$ is uniformly bounded from above by the area of the
round sphere of scalar curvature $\bar{R}$ for all large $a$. Thus
we deduce that $\bar{R}=1$. So \begin{equation}
{\rm Area}\,\{f=a\}<8\pi  %\eqno(6.4.8)
\end{equation} for $a$ large enough.

Denote by $X$ the unit normal vector to the level surface
$\{f=a\}$. By using the Gauss equation and (6.4.1), the intrinsic
curvature of the level surface $\{f=a\}$ can be computed as
\begin{align}
&\text{intrinsic curvature} \\
&  = R_{1212}+\det(h_{ij})  \nonumber\\
&  = \frac{1}{2}(R-2\Ric(X,X))
+\frac{\det({\rm Hess}\,(f))}{|\nabla f|^2} \nonumber\\
&  \le \frac{1}{2}(R-2\Ric(X,X))
+\frac{1}{4|\nabla f|^2}(\operatorname{tr}\,({\rm Hess}\,(f)))^2 \nonumber\\
&  = \frac{1}{2}(R-2\Ric(X,X))
+\frac{1}{4|\nabla f|^2}(1-(R-\Ric(X,X)))^2 \nonumber\\
& =\frac{1}{2}\left[1-\Ric(X,X)-(1-R+\Ric(X,X))
+\frac{(1-R+\Ric(X,X))^2}{2|\nabla f|^2}\right] \nonumber\\
&  < \frac{1}{2} \nonumber
\end{align} %\eqno(6.4.9)
for sufficiently large $a$, since $(1-R+Ric(X,X))>0$ and $|\nabla
f|$ is large when $a$ is large. Thus the combination of (6.4.8)
and (6.4.9) gives a contradiction to the Gauss-Bonnet formula.

Therefore we have proved the lemma.
%$$\eqno \#$$\vskip 0.3cm
\end{proof}

As a direct consequence, there is a universal positive constant
$\kappa_0$ such that any nonflat three-dimensional gradient
shrinking soliton, which is also an ancient $\kappa$-solution, to
the Ricci flow must be $\kappa_0$-noncollapsed on all scales
unless it is a metric quotient of  round three-sphere. The
following result, claimed by Perelman in the section 1.5 of
\cite{P2}, shows that this property actually holds for all nonflat
three-dimensional ancient $\kappa$-solutions.

%\vskip 0.2cm\noindent {\bf Proposition 6.4.2} (\textbf{
%CZ-SOURCE-LINE-15567
\begin{proposition}[universal noncollapsing of ancient three-dimensional solutions]
  \label{prop:cz-ancient-three-dimensional-universal-noncollapsing}
  \dcref{Cao--Zhu Proposition 6.4.2}
  \group{Three-Dimensional Ancient Kappa-Solutions}
  \source{cao-zhu}
 There exists a positive constant $\kappa_0$ with
the following property. Suppose we have a nonflat
three-dimensional ancient $\kappa$-solution for some $\kappa >0$.
Then either the solution is $\kappa_0$-noncollapsed on all scales,
or it is a metric quotient of the round three-sphere.
\end{proposition}

%\noindent{\textbf{Proof}.}
\begin{proof}
  \uses{lem:cz-three-dimensional-shrinking-soliton-classification, thm:cz-ancient-kappa-asymptotic-shrinker, lem:cz-unbounded-curvature-point-selection, prop:cz-pointed-limit-at-infinity-splits-line, prop:cz-neck-radius-lower-bound-at-infinity, prop:cz-three-dimensional-positive-ricci-asymptotic-roundness}
Let $g_{ij}(x,t), x\in M$ and $t\in(-\infty, 0]$, be a nonflat
ancient $\kappa$-solution for some $\kappa>0$. For an arbitrary
point $(p,t_0)\in M\times (-\infty,0]$, we define as in Chapter 3
that
\begin{align*}
& \tau=t_0-t,\quad \mbox{for }\;t<t_0,\\
& l(q,\tau)=\frac{1}{2\sqrt{\tau}}
\inf\bigg\{\int_0^\tau\sqrt{s}(R(\gamma(s),t_0-s)
+|\dot{\gamma}(s)|^2_{g_{ij}(t_0-s)})ds|\\
& \qquad \gamma:[0,\tau]\rightarrow M\;\mbox{ with }
\;\gamma(0)=p,\gamma(\tau)=q\bigg\},\\
&  \mbox{and }\widetilde{V}(\tau)
=\int_M(4\pi\tau)^{-\frac{3}{2}}\exp(-l(q,\tau))dV_{t_0-\tau}(q).
\end{align*}
Recall from (6.2.1) that for each $\tau>0$ we can find $q=q(\tau)$
such that $l(q,\tau)\le\frac{3}{2}$. In view of Lemma 6.4.1, we
may assume that the ancient $\kappa$-solution is not a gradient
shrinking Ricci soliton. Thus by (the proof of) Theorem 6.2.1, the
scalings of $g_{ij}(t_0-\tau)$ at $q(\tau)$ with factor
$\tau^{-1}$ converge along a subsequence of
$\tau\rightarrow+\infty$ to a nonflat gradient shrinking soliton
with nonnegative curvature operator which is $\kappa$-noncollapsed
on all scales. We now show that the limit has bounded curvature.

Denote the limiting nonflat gradient shrinking soliton by
$(\bar{M},\bar{g}_{ij}(x,t))$ with $-\infty < t \leq 0$. Note that
there holds the Li-Yau-Hamilton inequality (Theorem 2.5.4) on any
ancient $\kappa$-solution and in particular, the scalar curvature
of the ancient $\kappa$-solution is pointwise nondecreasing in
time. This implies that the scalar curvature of the limiting
soliton $(\bar{M},\bar{g}_{ij}(x,t))$ is still pointwise
nondecreasing in time. Thus we only need to show that the limiting
soliton has bounded curvature at $t=0$.

We argue by contradiction. By lifting to its orientable cover, we
may assume that $\bar{M}$ is orientable. Suppose the curvature of
the limiting soliton is unbounded at $t=0$. Of course in this case
the limiting soliton $\bar{M}$ is noncompact. Then by applying
Lemma 6.1.4, we can choose a sequence of points $x_j, j = 1, 2,
\ldots,$ divergent to infinity such that the scalar curvature
$\bar{R}$ of the limit satisfies
$$
\bar{R}(x_j,0) \geq j\; \mbox{ and }\; \bar{R}(x,0) \leq
4\bar{R}(x_j,0)
$$
for all $x\in B_{0}(x_j,{j}/{\sqrt{\bar{R}(x_{j},0)}})$ and $j =
1, 2,\, \ldots$.  Since the scalar curvature is nondecreasing in
time, we have \begin{equation}
\bar{R}(x,t) \leq 4\bar{R}(x_j,0), %\eqno(6.4.10)
\end{equation} for all $x\in B_{0}(x_j,{j}/{\sqrt{\bar{R}(x_{j},0)}})$, all
$t\leq 0$ and $j = 1, 2,\, \ldots$. By combining with Hamilton's
compactness theorem (Theorem 4.1.5) and the
$\kappa$-noncollapsing, we know that a subsequence of the rescaled
solutions
$$
(\bar{M},\bar{R}(x_j,0)\bar{g}_{ij}(x, t/\bar{R}(x_j,0)),x_j), \ \
j=1, 2, \ldots,
$$
converges in the $C^{\infty}_{\rm loc}$ topology to a nonflat
smooth solution of the Ricci flow. Then Proposition 6.1.2 implies
that the new limit at the new time $\{t = 0\}$ must split off a
line. By pulling back the new limit to its universal cover and
applying Hamilton's strong maximum principle, we deduce that the
pull-back of the new limit on the universal cover splits off a
line for all time $t \leq 0 $. Thus by combining with Theorem
6.2.2 and the argument in the proof of Lemma 6.4.1, we further
deduce that the new limit is either the round cylinder
$\mathbb{S}^2 \times \mathbb{R}$ or the round
$\mathbb{R}\mathbb{P}^2 \times \mathbb{R}$. Since $\bar{M}$ is
orientable, the new limit must be $\mathbb{S}^2 \times
\mathbb{R}$. Since $(\bar{M},\bar{g}_{ij}(x,0))$ has nonnegative
curvature operator and the points $\{ x_j \}$ going to infinity
and $ \bar{R}(x_j,0)\rightarrow +\infty $, this gives a
contradiction to Proposition 6.1.1. So we have proved that the
limiting gradient shrinking soliton has bounded curvature at each
time.

Hence by Lemma 6.4.1, the limiting gradient shrinking soliton is
either the round three-sphere $\mathbb{S}^3$ or its metric
quotients, or the infinite cylinder $\mathbb{S}^2\times
\mathbb{R}$ or one of its $\mathbb{Z}_2$ quotients. If the
asymptotic gradient shrinking soliton is the round three-sphere
$\mathbb{S}^3$ or its metric quotients, it follows from Lemma
5.2.4 and Proposition 5.2.5 that the ancient $\kappa$-solution
must be round. Thus in the following we may assume the asymptotic
gradient shrinking soliton is the infinite cylinder
$\mathbb{S}^2\times\mathbb{R}$ or a $\mathbb{Z}_2$ quotient of
$\mathbb{S}^2\times\mathbb{R}$.

We now come back to consider the original ancient
$\kappa$-solution $(M,g_{ij}(x,$ $t))$. By rescaling, we can
assume that $R(x,t)\le 1$ for all $(x,t)$ satisfying
$d_{t_0}(x,p)\le 2$ and $t\in[t_0-1,t_0]$. We will argue as in the
proof of Theorem 3.3.2 (Perelman's no local collapsing theorem I)
to obtain a positive lower bound for $\Vol_{t_0}(B_{t_0}(p,1))$.

Denote by $\xi=\Vol_{t_0}(B_{t_0}(p,1))^\frac{1}{3}$. For any
$v\in T_pM$ we can find an $\mathcal{L}$-geodesic $\gamma(\tau)$,
starting at $p$, with $\lim_{\tau\rightarrow
0^+}\sqrt{\tau}\dot{\gamma}(\tau)=v.$ It follows from the
$\mathcal{L}$-geodesic equation (3.2.1) that
$$
\frac{d}{d\tau}(\sqrt{\tau}\dot{\gamma})
-\frac{1}{2}\sqrt{\tau}\nabla
R+2\Ric(\sqrt{\tau}\dot{\gamma},\cdot)=0.
$$
By integrating as before we see that for $\tau\le\xi$ with the
property $\gamma(\sigma)\in B_{t_0}(p,1)$ as long as
$\sigma<\tau$, there holds
$$
|\sqrt{\tau}\dot{\gamma}(\tau)-v|\le C\xi (|v|+1)
$$
where $C$ is some positive constant depending only on the
dimension. Without loss of generality, we may assume
$C\xi\le\frac{1}{4}$ and $\xi\le\frac{1}{100}$. Then for $v\in
T_pM$ with $|v|\le\frac{1}{4}\xi^{-\frac{1}{2}}$ and for
$\tau\le\xi$ with the property $\gamma(\sigma)\in B_{t_0}(p,1)$ as
long as $\sigma<\tau$, we have
\begin{align*}
d_{t_0}(p,\gamma(\tau))
&  \le \int_0^\tau|\dot{\gamma}(\sigma)|d\sigma\\
&  < \frac{1}{2}\xi^{-\frac{1}{2}}
\int_0^\tau\frac{d\sigma}{\sqrt{\sigma}}\\
&  = 1.
\end{align*}
This shows \begin{equation}
\mathcal{L}\exp\left\{|v|\le\frac{1}{4}\xi^{-\frac{1}{2}}\right\}(\xi)
\subset B_{t_0}(p,1). %\eqno(6.4.11)
\end{equation}

We decompose Perelman's reduced volume $\widetilde{V}(\xi)$ as
\begin{align}
\widetilde{V}(\xi) &=\int_{\mathcal{L}\exp\left\{|v|
\le\frac{1}{4}\xi^{-\frac{1}{2}}\right\}(\xi)} \\
&\quad+\int_{M\setminus\mathcal{L}\exp\left\{|v|
\leq\frac{1}{4}\xi^{-\frac{1}{2}}\right\}(\xi)}(4\pi\xi)^{-\frac{3}{2}}
\exp(-l(q,\xi))dV_{t_0-\xi}(q). \nonumber%\eqno(6.4.12)
\end{align}
By using (6.4.11) and the metric evolution equation of the Ricci
flow, the first term on the RHS of (6.4.12) can be estimated by
\begin{align*}
&
\int_{\mathcal{L}\exp\{|v|\le\frac{1}{4}\xi^{-\frac{1}{2}}\}(\xi)}
(4\pi\xi)^{-\frac{3}{2}}\exp(-l(q,\xi))dV_{t_0-\xi}(q)\\
&  \le\int_{B_{t_0}(p,1)}(4\pi\xi)^{-\frac{3}{2}}e^{3\xi}dV_{t_0}(q)\\
&  = (4\pi)^{-\frac{3}{2}}e^{3\xi}{\xi}^{\frac{3}{2}}\\
&  < \xi^\frac{3}{2},
\end{align*}
while by using Theorem 3.2.7 (Perelman's Jacobian comparison
theorem), the second term on the RHS of (6.4.12) can be estimated
by
\begin{align}
& \int_{M\setminus\mathcal{L}\exp
\left\{|v|\leq\frac{1}{4}\xi^{-\frac{1}{2}}\right\}(\xi)}
(4\pi\xi)^{-\frac{3}{2}}
\exp(-l(q,\xi))dV_{t_0-\xi}(q)\\
& \le\int_{\{|v|>\frac{1}{4}\xi^{-\frac{1}{2}}\}}
(4\pi\tau)^{-\frac{3}{2}}
\exp(-l(\tau))\mathcal{J}(\tau)|_{\tau=0}dv \nonumber\\
&  =
(4\pi)^{-\frac{3}{2}}\int_{\{|v|>\frac{1}{4}\xi^{-\frac{1}{2}}\}}
\exp(-|v|^2)dv \nonumber\\
&  < \xi^\frac{3}{2} \nonumber
\end{align}  %\eqno(6.4.13)
since $\lim_{\tau\rightarrow
0^+}\tau^{-\frac{3}{2}}\mathcal{J}(\tau)=1$ and
$\lim_{\tau\rightarrow 0^+}l(\tau)=|v|^2$ by (3.2.18) and (3.2.19)
respectively. Thus we obtain \begin{equation}
\widetilde{V}(\xi)<2\xi^\frac{3}{2}. %\eqno(6.4.14)
\end{equation}

On the other hand, we recall that there exist a sequence
$\tau_k\rightarrow+\infty$ and a sequence of points $q(\tau_k)\in
M$ with $l(q(\tau_k),\tau_k)\le\frac{3}{2}$ so that the scalings
of the ancient $\kappa$-solution at $q(\tau_k)$ with factor
$\tau_k^{-1}$ converge to either round $\mathbb{S}^2\times
\mathbb{R}$ or one of its $\mathbb{Z}_2$ quotients. For
sufficiently large $k$, we construct a path
$\gamma:\quad[0,2\tau_k]\rightarrow M$, connecting $p$ to any
given point $q\in M$, as follows: the first half path
$\gamma|_{[0,\tau_k]}$ connects $p$ to $q(\tau_k)$ such that
$$
l(q(\tau_k),\tau_k)
=\frac{1}{2\sqrt{\tau_k}}\int_0^{\tau_k}\sqrt{\tau}
(R+|\dot{\gamma}(\tau)|^2)d\tau\le2,
$$
and the second half path $\gamma|_{[\tau_k,2\tau_k]}$ is a
shortest geodesic connecting $q(\tau_k)$ to $q$ with respect to
the metric $g_{ij}(t_0-\tau_k)$. Note that the rescaled metric
$\tau_k^{-1}g_{ij}(t_0-\tau)$ over the domain
$B_{t_0-\tau_k}(q(\tau_k),\sqrt{\tau_k})\times[t_0-2\tau_k,t_0-\tau_k]$
is sufficiently close to the round $\mathbb{S}^2\times \mathbb{R}$
or its $\mathbb{Z}_2$ quotients. Then there is a universal
positive constant $\beta$ such that
\begin{align*}
l(q,2\tau_k)& \le \frac{1}{2\sqrt{2\tau_k}}\left(\int_0^{\tau_k}
+\int_{\tau_k}^{2\tau_k}\right)\sqrt{\tau}
(R+|\dot{\gamma}(\tau)|^2)d\tau\\
&  \le \sqrt{2}+\frac{1}{2\sqrt{2\tau_k}}
\int_{\tau_k}^{2\tau_k}\sqrt{\tau}(R+|\dot{\gamma}(\tau)|^2)d\tau\\
&  \le \beta
\end{align*}
for all $q\in B_{t_0-\tau_k}(q(\tau_k),\sqrt{\tau_k})$. Thus
\begin{align*}
\widetilde{V}(2\tau_k)&  = \int_M (4\pi(2\tau_k))^{-\frac{3}{2}}
\exp(-l(q,2\tau_k))dV_{t_0-2\tau_k}(q)\\
&  \ge e^{-\beta}\int_{B_{t_0-\tau_k}(q(\tau_k),
\sqrt{\tau_k})}(4\pi(2\tau_k))^{-\frac{3}{2}}dV_{t_0-2\tau_k}(q)\\
&  \ge \tilde{\beta}
\end{align*}
for some universal positive constant $\tilde{\beta}$. Here we have
used the curvature estimate (6.2.6). By combining with the
monotonicity of Perelman's reduced volume (Theorem 3.2.8) and
(6.4.14), we deduce that
$$
\tilde{\beta}\le\widetilde{V}(2\tau_k)
\le\widetilde{V}(\xi)<2\xi^\frac{3}{2}.
$$
This proves
$$
\Vol_{t_0}(B_{t_0}(p,1))\ge \kappa_0>0
$$
for some universal positive constant $\kappa_0$. So we have proved
that the ancient $\kappa$-solution is also an ancient
$\kappa_0$-solution.
%$$\eqno \#$$ \vskip 0.3cm
\end{proof}

The important Li-Yau-Hamilton inequality gives rise to a parabolic
Harnack estimate (Corollary 2.5.7) for solutions of the Ricci flow
with bounded and nonnegative curvature operator. As explained in
the previous section, the no local collapsing theorem of Perelman
implies a volume lower bound from a curvature upper bound, while
the estimate in the previous section implies a curvature upper
bound from a volume lower bound. The combination of these two
estimates as well as the Li-Yau-Hamilton inequality will give an
important elliptic type property for three-dimensional ancient
$\kappa$-solutions. This elliptic type property was first
implicitly given by Perelman in \cite{P1} and it will play a
crucial role in the analysis of singularities.

%\vskip 0.2cm\noindent {\bf Theorem 6.4.3} (\textbf{
%CZ-SOURCE-LINE-15819
\begin{theorem}[elliptic estimates for ancient kappa-solutions]
  \label{thm:cz-ancient-kappa-elliptic-estimates}
  \dcref{Cao--Zhu Theorem 6.4.3}
  \group{Three-Dimensional Ancient Kappa-Solutions}
  \source{cao-zhu}
 There exist a positive constant $\eta$ and a
positive increasing function $\omega:\quad [0,+\infty)\rightarrow
(0,+\infty)$ with the following properties. Suppose we have a
three-dimensional ancient $\kappa$-solution
$(M,g_{ij}(t)),-\infty<t\le 0,$ for some $\kappa>0$. Then
\begin{itemize}
\item[(i)]for every $x,y\in M$ and $t\in(-\infty,0]$, there holds
$$
R(x,t)\le R(y,t)\cdot \omega(R(y,t)d_t^2(x,y));
$$
\item[(ii)] for all $x\in M$ and $t\in (-\infty,0]$, there hold
$$
|\nabla R|(x,t)\le\eta R^\frac{3}{2}(x,t)\;  \mbox{ and }\;
\left|\frac{\partial R}{\partial t}\right|(x,t)\le \eta R^2(x,t).
$$
\end{itemize}
\end{theorem}

%\vskip0.1cm \noindent{\textbf{Proof.}}
\begin{proof}
  \uses{thm:cz-local-volume-growth-curvature-estimate} \

(i) Consider a three-dimensional nonflat ancient $\kappa$-solution
$g_{ij}(x,t)$ on $M \times (-\infty,0]$. In view of Proposition
6.4.2, we may assume that the ancient solution is universal
$\kappa_0$-noncollapsed. Obviously we only need to establish the
estimate at $t=0$. Let $y$ be an arbitrarily fixed point in $M$.
By rescaling, we can assume $R(y,0)=1$.

Let us first consider the case that $\sup\{R(x,0)d_0^2(x,y)\ |\
x\in M\}>1$. Define $z$ to be the closest point to $y$ (at time
$t=0$) satisfying $R(z,0)d_0^2(z,y)=1$. We want to bound
$R(x,0)/R(z,0)$ from above for $x\in
B_0(z,2R(z,0)^{-\frac{1}{2}}).$

Connect $y$ and $z$ by a shortest geodesic and choose a point
$\tilde{z}$ lying on the geodesic satisfying
$d_0(\tilde{z},z)=\frac{1}{4}R(z,0)^{-\frac{1}{2}}$. Denote by $B$
the ball centered at $\tilde{z}$ and with radius
$\frac{1}{4}R(z,0)^{-\frac{1}{2}}$ (with respect to the metric at
$t=0$). Clearly the ball $B$ lies in
$B_0(y,R(z,0)^{-\frac{1}{2}})$ and lies outside
$B_0(y,\frac{1}{2}R(z,0)^{-\frac{1}{2}})$. Thus for $x\in B$, we
have
$$
R(x,0)d_0^2(x,y)\le 1\quad \mbox{and} \quad
d_0(x,y)\ge\frac{1}{2}R(z,0)^{-\frac{1}{2}}
$$
and hence
$$
R(x,0)\le\frac{1}{(\frac{1}{2}R(z,0)^{-\frac{1}{2}})^2} \quad
\mbox{for all } \;x\in B.
$$
Then by the Li-Yau-Hamilton inequality and the
$\kappa_0$-noncollapsing, we have
$$
\Vol_0(B)\ge\kappa_0\left(\frac{1}{4}R(z,0)^{-\frac{1}{2}}\right)^3,
$$
and then
$$
\Vol_0(B_0(z,8R(z,0)^{-\frac{1}{2}}))
\ge\frac{\kappa_0}{2^{15}}(8R(z,0)^{-\frac{1}{2}})^3.
$$
So by Theorem 6.3.3(ii), there exist positive constants
$B(\kappa_0), C(\kappa_0),$ and $\tau_0(\kappa_0)$ such that \begin{equation}
R(x,0)\le (C(\kappa_0)+\frac{B(\kappa_0)}{\tau_0(\kappa_0)})R(z,0)
%\eqno(6.4.15)
\end{equation} for all $x\in B_0(z,2R(z,0)^{-\frac{1}{2}})$.

We now consider the remaining case. If $R(x,0)d_0^2(x,y)\le 1$
everywhere, we choose a point $z$ satisfying $\sup\{R(x,0)\ |\
x\in M\}\leq2R(z,0)$. Obviously we also have the estimate (6.4.15)
in this case.

We next want to bound $R(z,0)$ for the chosen $z\in M$. By
(6.4.15) and the Li-Yau-Hamilton inequality, we have
$$
R(x,t)\le (C(\kappa_0)+\frac{B(\kappa_0)}{\tau_0(\kappa_0)})R(z,0)
$$
for all $x\in B_0(z,2R(z,0)^{-\frac{1}{2}})$ and all $t\le 0$. It
then follows from the local derivative estimates of Shi that
$$
\frac{\partial R}{\partial t}(z,t)
\le\widetilde{C}(\kappa_0)R(z,0)^2,\quad \mbox{for all }\;
-R^{-1}(z,0)\leq t\leq 0
$$
which implies \begin{equation}
R(z,-cR^{-1}(z,0))\ge cR(z,0)  %\eqno(6.4.16)
\end{equation} for some small positive constant $c$ depending only on
$\kappa_0$. On the other hand, by using the Harnack estimate in
Corollary 2.5.7, we have \begin{equation}
1=R(y,0)\ge\widetilde{c}R(z,-cR^{-1}(z,0)) %\eqno(6.4.17)
\end{equation} for some small positive constant $\widetilde{c}$ depending
only on $\kappa_0$, since $d_0(y,z)\le R(z,0)^{-\frac{1}{2}}$ and
the metric $g_{ij}(t)$ is equivalent on
$$
B_0(z,2R(z,0)^{-\frac{1}{2}})\times[-cR^{-1}(z,0),0]
$$
with $c>0$ small enough. Thus we get from (6.4.16) and (6.4.17)
that \begin{equation}
R(z,0)\le\widetilde{A} %\eqno(6.4.18)
\end{equation} for some positive constant $\widetilde{A}$ depending only on
$\kappa_0$.

Since $B_0(z,2R(z,0)^{-\frac{1}{2}})\supset
B_0(y,R(z,0)^{-\frac{1}{2}})$ and $R(z,0)^{-\frac{1}{2}}\ge
(\widetilde{A})^{-\frac{1}{2}}$, the combination of (6.4.15) and
(6.4.18) gives \begin{equation} R(x,0)\le
(C(\kappa_0)+\frac{B(\kappa_0)}{\tau_0(\kappa_0)})
\widetilde{A}  %\eqno(6.4.19)
\end{equation} whenever $x\in B_0(y,(\widetilde{A})^{-\frac{1}{2}})$. Then by
the $\kappa_0$-noncollapsing there exists a positive constant
$r_0$ depending only on $\kappa_0$ such that
$$
\Vol_0(B_0(y,r_0))\ge \kappa_0r_0^3.
$$
For any fixed $R_0\ge r_0$, we then have
$$
\Vol_0(B_0(y,R_0))\ge \kappa_0r_0^3
=\kappa_0(\frac{r_0}{R_0})^3\cdot R_0^3.
$$
By applying Theorem 6.3.3 (ii) again and noting that the constant
$\kappa_0$ is universal, there exists a positive constant
$\omega(R_0)$ depending only on $R_0$ such that
$$
R(x,0)\le\omega(R_0^2)\qquad \mbox{for all }\; x\in
B_0(y,\frac{1}{4}R_0).
$$
This gives the desired estimate.

\medskip
(ii) This follows immediately from  conclusion (i), the
Li-Yau-Hamilton inequality and the local derivative estimate of
Shi.
%$$\eqno \#$$ \vskip 0.3cm
\end{proof}

As a consequence, we have the following compactness result due to
Perelman \cite{P1}.

%\vskip 0.2cm\noindent{\bf Corollary 6.4.4} (\textbf{
%CZ-SOURCE-LINE-15960
\begin{corollary}[compactness of normalized ancient kappa-solutions]
  \label{cor:cz-ancient-kappa-solution-compactness}
  \dcref{Cao--Zhu Corollary 6.4.4}
  \group{Three-Dimensional Ancient Kappa-Solutions}
  \source{cao-zhu}
 The set of
nonflat three-dimensional ancient $\kappa_0$-solutions is compact
modulo scaling in the sense that for any sequence of such
solutions and marking points $(x_k,0)$ with $R(x_k,0)=1$, we can
extract a $C^\infty_{loc}$ converging subsequence whose limit is
also an ancient $\kappa_0$-solution.
\end{corollary}

%\vskip0.1cm \noindent{\textbf{Proof}.}
\begin{proof}
  \uses{thm:cz-ancient-kappa-elliptic-estimates, lem:cz-unbounded-curvature-point-selection, prop:cz-pointed-limit-at-infinity-splits-line, thm:cz-two-dimensional-ancient-kappa-classification, lem:cz-three-dimensional-shrinking-soliton-classification, prop:cz-neck-radius-lower-bound-at-infinity}
Consider any sequence of three-dimensional ancient
$\kappa_0$-solutions and marking points $(x_k,0)$ with
$R(x_k,0)=1$. By Theorem 6.4.3(i), the Li-Yau-Hamilton inequality
and Hamilton's compactness theorem (Theorem 4.1.5), we can extract
a $C^\infty_{loc}$ converging subsequence such that the limit
$(\bar{M},\bar{g}_{ij}(x,t))$, with $-\infty <t \leq 0$, is an
ancient solution to the Ricci flow with nonnegative curvature
operator and $\kappa_0$-noncollapsed on all scales. Since any
ancient $\kappa_0$-solution satisfies the Li-Yau-Hamilton
inequality, it implies that the scalar curvature $\bar{R}(x,t)$ of
the limit $(\bar{M},\bar{g}_{ij}(x,t))$ is pointwise nondecreasing
in time. Thus it remains to show that the limit solution has
bounded curvature at $t=0$.

Obviously we may assume the limiting manifold $\bar{M}$ is
noncompact. By pulling back the limiting solution to its
orientable cover, we can assume that the limiting manifold
$\bar{M}$ is orientable. We now argue by contradiction. Suppose
the scalar curvature $\bar{R}$ of the limit at $t=0$ is unbounded.

By applying Lemma 6.1.4, we can choose a sequence of points $x_j
\in \bar{M}, j = 1, 2, \ldots,$ divergent to infinity such that
the scalar curvature $\bar{R}$ of the limit satisfies
$$
\bar{R}(x_j,0) \geq j\; \mbox{ and }\; \bar{R}(x,0) \leq
4\bar{R}(x_j,0)
$$
for all $j = 1, 2,\, \ldots,$ and $x\in
B_{0}(x_j,{j}/{\sqrt{\bar{R}(x_{j},0)}})$. Then from the fact that
the limiting scalar curvature $\bar{R}(x,t)$ is pointwise
nondecreasing in time, we have \begin{equation}
\bar{R}(x,t) \leq 4\bar{R}(x_j,0) %\eqno(6.4.20)
\end{equation} for all $j = 1, 2,\, \ldots$, $x\in
B_{0}(x_j,{j}/{\sqrt{\bar{R}(x_{j},0)}})$ and $t\leq 0$. By
combining with Hamilton's compactness theorem (Theorem 4.1.5) and
the $\kappa_0$-noncollapsing, we know that a subsequence of the
rescaled solutions
$$
(\bar{M},\bar{R}(x_j,0)\bar{g}_{ij}(x, t/\bar{R}(x_j,0)),x_j), \ \
j=1, 2, \ldots,
$$
converges in the $C^{\infty}_{loc}$ topology to a nonflat smooth
solution of the Ricci flow. Then Proposition 6.1.2 implies that
the new limit at the new time $\{t = 0\}$ must split off a line.
By pulling back the new limit to its universal cover and applying
Hamilton's strong maximum principle, we deduce that the pull-back
of the new limit on the universal cover splits off a line for all
time $t \leq 0$. Thus by combining with Theorem 6.2.2 and the
argument in the proof of Lemma 6.4.1, we further deduce that the
new limit is either the round cylinder $\mathbb{S}^2 \times
\mathbb{R}$ or the round $\mathbb{R}\mathbb{P}^2 \times
\mathbb{R}$. Since $\bar{M}$ is orientable, the new limit must be
$\mathbb{S}^2 \times \mathbb{R}$. Moreover, since
$(\bar{M},\bar{g}_{ij}(x,0))$ has nonnegative curvature operator
and the points $\{ x_j \}$ are going to infinity and $
\bar{R}(x_j,0)\rightarrow +\infty $, this gives a contradiction to
Proposition 6.1.1. So we have proved that the limit
$(\bar{M},\bar{g}_{ij}(x,t))$ has uniformly bounded curvature.
%$$\eqno \#$$ \vskip 0.3cm
\end{proof}

Arbitrarily fix $\varepsilon>0$. Let $g_{ij}(x,t)$ be a nonflat
ancient $\kappa$-solution on a three-manifold $M$ for some $\kappa
>0$. We say that a point $x_0\in M$ is the {\bf center of an
evolving $\varepsilon$-neck}\index{center of an evolving
$\varepsilon$-neck} at $t=0$, if the solution $g_{ij}(x,t)$ in the
set $\{(x,t)\ |\ -\varepsilon^{-2} Q^{-1}<t\le 0,
d_t^2(x,x_0)<\varepsilon^{-2} Q^{-1}\}$, where $Q=R(x_0,0),$ is,
after scaling with factor $Q$, $\varepsilon$-close (in the
$C^{[\varepsilon^{-1}]}$ topology) to the corresponding set of the
evolving round cylinder having scalar curvature one at $t=0$.

As another consequence of the elliptic type estimate, we have the
following global structure result obtained by Perelman in
\cite{P1} for noncompact ancient $\kappa$-solutions.

%\vskip 0.2cm\noindent{\bf Corollary 6.4.5}.  \emph{
%CZ-SOURCE-LINE-16048
\begin{corollary} [cap-or-neck structure of noncompact ancient solutions]
  \label{cor:cz-noncompact-ancient-cap-neck-structure}
  \dcref{Cao--Zhu Corollary 6.4.5}
  \group{Three-Dimensional Ancient Kappa-Solutions}
  \source{cao-zhu}

For any $\varepsilon>0$ there exists $C=C(\varepsilon)>0$, such
that if $g_{ij}(t)$ is a nonflat ancient $\kappa$-solution on a
noncompact three-manifold $M$ for some $\kappa>0$, and
$M_\varepsilon$ denotes the set of points in $M$ which are not
centers of evolving $\varepsilon$-necks at $t=0$, then at $t=0$,
either the whole manifold $M$ is the round cylinder $\mathbb{S}^2
\times \mathbb{R}$ or its $\mathbb{Z}_2$ metric quotients, or
$M_\varepsilon$ satisfies the following
\begin{itemize}
\item[(i)] $M_\varepsilon$ is compact, \item[(ii)] ${\rm diam}\,
M_\varepsilon \le CQ^{-\frac{1}{2}}$ and $C^{-1}Q\le R(x,0)\le
CQ$, whenever $x\in M_\varepsilon$, where $Q=R(x_0,0)$ for some
$x_0\in\partial M_\varepsilon$.
\end{itemize}
\end{corollary}

%\vskip0.1cm \noindent{\textbf{Proof.}}
\begin{proof}
  \uses{lem:cz-three-dimensional-shrinking-soliton-classification, prop:cz-pointed-limit-at-infinity-splits-line, thm:cz-two-dimensional-ancient-kappa-classification, cor:cz-ancient-kappa-solution-compactness, prop:cz-ancient-three-dimensional-universal-noncollapsing, thm:cz-ancient-kappa-elliptic-estimates}
We first consider the easy case that the curvature operator of the
ancient $\kappa$-solution has a nontrivial null vector somewhere
at some time. Let us pull back the solution to its universal
cover. By applying Hamilton's strong maximum principle and Theorem
6.2.2, we see that the universal cover is the evolving round
cylinder $\mathbb{S}^2 \times \mathbb{R}$. Thus in this case, by
the argument in the proof of Lemma 6.4.1, we conclude that the
ancient $\kappa$-solution is either isometric to the round
cylinder $\mathbb{S}^2 \times \mathbb{R}$ or one of its
$\mathbb{Z}_2$ metric quotients (i.e., $\mathbb{R}\mathbb{P}^2
\times \mathbb{R}$, or the twisted product $\mathbb{S}^2
\tilde{\times} \mathbb{R}$ where $\mathbb{Z}_2$ flips both
$\mathbb{S}^2$, or $\mathbb{R}$).

We then assume that the curvature operator of the nonflat ancient
$\kappa$-solution is positive everywhere. Firstly we want to show
$M_\varepsilon$ is compact. We argue by contradiction. Suppose
there exists a sequence of points $z_k,\ k=1,2,\ldots$, going to
infinity (with respect to the metric $g_{ij}(0)$) such that each
$z_k$ is not the center of any evolving $\varepsilon$-neck. For an
arbitrarily fixed point $z_0\in M$, it follows from Theorem
6.4.3(i) that
$$
0<R(z_0,0)\leq R(z_k,0)\cdot \omega(R(z_k,0)d^2_0(z_k,z_0))
$$
which implies that
$$
\lim_{k\rightarrow\infty}R(z_k,0)d^2_0(z_k,z_0)=+\infty.
$$
Since the sectional curvature of the ancient $\kappa$-solution is
positive everywhere, the underlying manifold is diffeomorphic to
$\mathbb{R}^3$, and in particular, orientable. Then as before, by
Proposition 6.1.2, Theorem 6.2.2 and Corollary 6.4.4, we conclude
that $z_k$ is the center of an evolving $\varepsilon$-neck for $k$
sufficiently large. This is a contradiction, so we have proved
that $M_{\varepsilon}$ is compact.

Again, we notice that $M$ is diffeomorphic to $\mathbb{R}^3$ since
the curvature operator is positive. According to the resolution of
the Schoenflies conjecture in three-dimensions, every
approximately round two-sphere cross-section through the center of
an evolving $\varepsilon$-neck divides $M$ into two parts such
that one of them is diffeomorphic to the three-ball
$\mathbb{B}^3$. Let $\varphi$ be the Busemann function on $M$, it
is a standard fact that $\varphi$ is convex and proper. Since
$M_{\varepsilon}$ is compact, $M_\varepsilon$ is contained in a
compact set $K=\varphi^{-1}((-\infty,A])$ for some large $A$. We
note that each point $x\in M\setminus M_\varepsilon$ is the center
of an $\varepsilon$-neck. It is clear that there is an
$\varepsilon$-neck $N$ lying entirely outside $K$.  Consider a
point $x$ on one of the boundary components of the
$\varepsilon$-neck $N$. Since $x \in M\setminus M_{\varepsilon}$,
there is an $\varepsilon$-neck adjacent to the initial
$\varepsilon$-neck, producing a longer neck. We then take a point
on the boundary of the second $\varepsilon$-neck and continue.
This procedure can either terminate when we get into
$M_{\varepsilon}$ or go on infinitely to produce a semi-infinite
(topological) cylinder. The same procedure can be repeated for the
other boundary component of the initial $\varepsilon$-neck. This
procedure will give a maximal extended neck $\tilde{N}$. If
$\tilde{N}$ never touches $M_\varepsilon$, the manifold will be
diffeomorphic to the standard infinite cylinder, which is a
contradiction. If both ends of $\tilde{N}$ touch $M_\varepsilon$,
then there is a geodesic connecting two points of
$M_{\varepsilon}$ and passing through $N$. This is impossible
since the function $\varphi$ is convex. So we conclude that one
end of $\tilde{N}$ will touch $M_\varepsilon$ and the other end
will tend to infinity to produce a semi-infinite (topological)
cylinder. Thus we can find an approximately round two-sphere
cross-section which encloses the whole set $M_{\varepsilon}$ and
touches some point $x_{0}\in
\partial M_{\varepsilon}$. We next want to show that
$R(x_0,0)^{\frac{1}{2}}\cdot diam(M_{\varepsilon})$ is bounded
from above by some positive constant $C=C(\varepsilon)$ depending
only on $\varepsilon$.

Suppose not; then there exists a sequence of nonflat noncompact
three-dimensional ancient $\kappa$-solutions with positive
curvature operator such that for the above chosen points $x_0\in
\partial M_{\varepsilon}$ there would hold
\begin{equation} R(x_0,0)^{\frac{1}{2}}\cdot
{\rm diam}\,(M_{\varepsilon})\rightarrow +\infty. %\eqno(6.4.21)
\end{equation} By Proposition 6.4.2, we know that the ancient solutions are
$\kappa_0$-noncollapsed on all scales for some universal positive
constant $\kappa_0$. Let us dilate the ancient solutions around
the points $x_0$ with the factors $R(x_0,0)$. By Corollary 6.4.4,
we can extract a convergent subsequence. From the choice of the
points $x_0$ and (6.4.21), the limit has at least two ends. Then
by Toponogov's splitting theorem the limit is isometric to
$X\times\mathbb{R} $ for some nonflat two-dimensional ancient
$\kappa_0$-solution $X$. Since $M$ is orientable, we conclude from
Theorem 6.2.2 that limit must be the evolving round cylinder
$\mathbb{S}^2\times\mathbb{R}$. This contradicts the fact that
each chosen point $x_0$ is not the center of any evolving
$\varepsilon$-neck. Therefore we have proved
$$
{\rm diam}\,(M_{\varepsilon})\leq CQ^{-\frac{1}{2}}
$$
for some positive constant $C=C(\varepsilon)$ depending only on
$\varepsilon$, where $Q=R(x_0,0)$.

Finally by combining this diameter estimate with Theorem 6.4.3(i),
we immediately deduce
$$
\widetilde{C}^{-1}Q\leq R(x,0)\leq \widetilde{C}Q,\;\mbox{
whenever }\;x\in M_{\varepsilon},
$$
for some positive constant $\widetilde{C}$ depending only on
$\varepsilon$.
%$$\eqno \#$$\vskip 0.3cm
\end{proof}

We now can describe the canonical structures for three-dimensional
nonflat (compact or noncompact) ancient $\kappa$-solutions. The
following theorem was given by Perelman in the section 1.5 of
\cite{P2}. Recently in \cite{CZ05F}, this canonical neighborhood
result has been extended to four-dimensional ancient
$\kappa$-solutions with isotropic curvature pinching.

%\vskip0.2cm\noindent{\bf Theorem 6.4.6} (\textbf
%CZ-SOURCE-LINE-16187
\begin{theorem}[canonical neighborhoods in ancient kappa-solutions]
  \label{thm:cz-ancient-kappa-canonical-neighborhoods}
  \dcref{Cao--Zhu Theorem 6.4.6}
  \group{Three-Dimensional Ancient Kappa-Solutions}
  \source{cao-zhu}
 For any $\varepsilon>0$ one can find
positive constants $C_1=C_1(\varepsilon)$ and
$C_2=C_2(\varepsilon)$ with the following property. Suppose we
have a three-dimensional nonflat $($compact or noncompact$)$
ancient $\kappa$-solution $(M,g_{ij}(x,t))$. Then either the
ancient solution is the round $\mathbb{R}\mathbb{P}^2 \times
\mathbb{R}$, or every point $(x,t)$ has an open neighborhood $B$,
with $B_t(x,r)\subset B \subset B_t(x,2r)$ for some
$0<r<C_1R(x,t)^{-\frac{1}{2}}$, which falls into one of the
following three categories:
\begin{itemize}
\item[(a)] $B$ is an {\bf evolving $\varepsilon$-neck} $($in the
sense that it is the slice at the time $t$ of the parabolic region
$\{(x',t')\ |\ x' \in B, t' \in [t - \varepsilon^{-2}R(x,t)^{-1},
t] \}$ which is, after scaling with factor $R(x,t)$ and shifting
the time $t$ to zero, $\varepsilon$-close $($in the
$C^{[\varepsilon^{-1}]}$ topology$)$ to the subset
$(\mathbb{S}^2\times \mathbb{I})\times [-\varepsilon^{-2},0]$ of
the evolving standard round cylinder with scalar curvature $1$ and
length $2\varepsilon^{-1}$ to $\mathbb{I}$ at the time zero$),$ or
\index{evolving $\varepsilon$-neck} \item[(b)] $B$ is an {\bf
evolving $\varepsilon$-cap} $($in the sense that it is the time
slice at the time $t$ of an evolving metric on $\mathbb{B}^3$ or
$\mathbb{R}\mathbb{P}^3\setminus\bar{\mathbb{B}^3}$ such that the
region outside some suitable compact subset of $\mathbb{B}^3$ or
$\mathbb{R}\mathbb{P}^3\setminus\bar{\mathbb{B}^3}$ is an evolving
$\varepsilon$-neck$),$ or \index{evolving $\varepsilon$-cap}
\item[(c)] $B$ is a compact manifold $($without boundary$)$ with
positive sectional curvature $($thus it is diffeomorphic to the
round three-sphere $\mathbb{S}^3$ or a metric quotient of
$\mathbb{S}^3);$
\end{itemize}
furthermore, the scalar curvature of the ancient $\kappa$-solution
on $B$ at time $t$ is between $C_2^{-1}R(x,t)$ and $C_2R(x,t)$,
and the volume of $B$ in case {\rm (a)} and case {\rm (b)}
satisfies
$$
(C_2R(x,t))^{-\frac{3}{2}} \leq \Vol_t(B) \leq \varepsilon r^3.
$$
\end{theorem}

%\vskip 0.1cm \noindent{\textbf{Proof.}}
\begin{proof}
  \uses{thm:cz-two-dimensional-ancient-kappa-classification, lem:cz-three-dimensional-shrinking-soliton-classification, cor:cz-noncompact-ancient-cap-neck-structure, thm:cz-ancient-kappa-elliptic-estimates, prop:cz-ancient-three-dimensional-universal-noncollapsing, cor:cz-ancient-kappa-solution-compactness}
As before, we first consider the easy case that the curvature
operator has a nontrivial null vector somewhere at some time. By
pulling back the solution to its universal cover and applying
Hamilton's strong maximum principle and Theorem 6.2.2, we deduce
that the universal cover is the evolving round cylinder
$\mathbb{S}^2 \times \mathbb{R}$. Then exactly as before, by the
argument in the proof of Lemma 6.4.1, we conclude that the ancient
$\kappa$-solution is isometric to the round $\mathbb{S}^2 \times
\mathbb{R}$, $\mathbb{R}\mathbb{P}^2 \times \mathbb{R}$, or the
twisted product $\mathbb{S}^2 \tilde{\times} \mathbb{R}$ where
$\mathbb{Z}_2$ flips both $\mathbb{S}^2$ and $\mathbb{R}$. Clearly
each point of the round cylinder $\mathbb{S}^2 \times \mathbb{R}$
or the twisted product $\mathbb{S}^2 \tilde{\times} \mathbb{R}$
has a neighborhood falling into the category (a) or (b) (over
$\mathbb{R}\mathbb{P}^3\setminus\bar{\mathbb{B}^3}$).

We now assume that the curvature operator of the nonflat ancient
$\kappa$-solution is positive everywhere. Then the manifold is
orientable by the Cheeger-Gromoll theorem \cite{CG} for the
noncompact case or the Synge theorem \cite{CE} for the compact
case.

Without loss of generality, we may assume $\varepsilon$ is
suitably small, say $0<\varepsilon < \frac{1}{100}$. If the
nonflat ancient $\kappa$-solution is noncompact, the conclusions
follow immediately from the combination of Corollary 6.4.5 and
Theorem 6.4.3(i). Thus we may assume the nonflat ancient
$\kappa$-solution is compact. By Proposition 6.4.2, either the
compact ancient $\kappa$-solution is isometric to a metric
quotient of the round $\mathbb{S}^3$, or it is
$\kappa_0$-noncollapsed on all scales for the universal positive
constant $\kappa_0$. Clearly each point of a metric quotient of
the round $\mathbb{S}^3$ has a neighborhood falling into category
(c). Thus we may further assume the ancient $\kappa$-solution is
also $\kappa_0$-noncollapsing.

We argue by contradiction. Suppose that for some $\varepsilon \in
(0, \frac{1}{100})$, there exist a sequence of compact orientable
ancient $\kappa_0$-solutions $(M_k,g_k)$ with positive curvature
operator, a sequence of points $(x_k,0)$ with $x_k\in M_k$ and
sequences of positive constants $C_{1k}\rightarrow\infty$ and
$C_{2k}=\omega(4C^2_{1k})$, with the function $\omega$ given in
Theorem 6.4.3, such that for every radius $r$,
$0<r<C_{1k}R(x_k,0)^{-\frac{1}{2}}$, any open neighborhood $B$,
with $B_0(x_k,r) \subset B \subset B_0(x_k,2r)$, does not fall
into one of the three categories (a), (b) and (c), where in the
case (a) and case (b), we require the neighborhood $B$ to satisfy
the volume estimate
$$
(C_{2k}R(x_k,0))^{-\frac{3}{2}} \leq \Vol_0(B) \leq \varepsilon
r^3.
$$

By Theorem 6.4.3(i) and the choice of the constants $C_{2k}$ we
see that the diameter of each $M_k$ at $t=0$ is at least
$C_{1k}R(x_k,0)^{-\frac{1}{2}}$; otherwise we can choose suitable
$r \in (0,C_{1k}R(x_k,0)^{-\frac{1}{2}})$ and $B=M_k$, which falls
into the category (c) with the scalar curvature between
$C_{2k}^{-1}R(x,0)$ and $C_{2k}R(x,0)$ on $B$. Now by scaling the
ancient $\kappa_0$-solutions along the points $(x_k,0)$ with
factors $R(x_k,0)$, it follows from Corollary 6.4.4 that a
sequence of the ancient $\kappa_0$-solutions converge in the
$C^{\infty}_{loc}$ topology to a noncompact orientable ancient
$\kappa_0$-solution.

If the curvature operator of the noncompact limit has a nontrivial
null vector somewhere at some time, it follows exactly as before
by using the argument in the proof of Lemma 6.4.1 that the
orientable limit is isometric to the round $\mathbb{S}^2 \times
\mathbb{R}$, or the twisted product $\mathbb{S}^2 \tilde{\times}
\mathbb{R}$ where $\mathbb{Z}_2$ flips both $\mathbb{S}^2$ and
$\mathbb{R}$. Then for $k$ large enough, a suitable neighborhood
$B$ (for suitable $r$) of the point $(x_k,0)$ would fall into the
category (a) or (b) (over
$\mathbb{R}\mathbb{P}^3\setminus\bar{\mathbb{B}^3}$) with the
desired volume estimate. This is a contradiction.

If the noncompact limit has positive sectional curvature
everywhere, then by using Corollary 6.4.5 and Theorem 6.4.3(i) for
the noncompact limit we see that for $k$ large enough, a suitable
neighborhood $B$ (for suitable $r$) of the point $(x_k,0)$ would
fall into category (a) or (b) (over $\mathbb{B}^3$) with the
desired volume estimate. This is also a contradiction.

Finally, the statement on the curvature estimate in the
neighborhood $B$ follows directly from Theorem 6.4.3(i).
%$$\eqno \#$$
\end{proof}


%\bigskip
\newpage

\section{Named intermediate formalization obligations}

\begin{lemma}[uniform reduced-distance and curvature bounds on parabolic annuli]
  \label{lem:cz-reduced-distance-curvature-parabolic-annulus}
  \dcref{Claim (6.2.6)}
  \group{Asymptotic Shrinking Solitons}
  \source{cao-zhu}
We first claim that for any $A \geq 1$, one can find
$B=B(A)<+\infty$ such that for every $\bar{\tau}>1$ there holds
\begin{equation} l(q,\tau)\leq B \; \mbox{ and }\;  \tau R(q,t_0-\tau)\leq B,
%\eqno(6.2.6)
\end{equation} whenever $\frac{1}{2}\bar{\tau}\leq \tau \leq A\bar{\tau}$ and
$d_{t_0-\frac{\bar{\tau}}{2}}^2 (q,q(\frac{\bar{\tau}}{2}))\leq
A\bar{\tau}.$

\end{lemma}
\begin{proof}
\begin{align}
\sqrt{l(q,\frac{\bar{\tau}}{2})} &  \leq \sqrt{\frac{n}{2}}+
\sup\{|\nabla \sqrt{l}|\}\cdot d_{t_0-\frac{\bar{\tau}}{2}}
\left(q,q\left(\frac{\bar{\tau}}{2}\right)\right) \\
&  \leq \sqrt{\frac{n}{2}}+ \sqrt{\frac{3}{2\bar{\tau}}}\cdot
\sqrt{A\bar{\tau}} \nonumber\\
&  = \sqrt{\frac{n}{2}}+\sqrt{\frac{3A}{2}},\nonumber
\end{align}  %\eqno(6.2.7)
and
\begin{align}
R\left(q,t_0-\frac{\bar{\tau}}{2}\right)&  \leq
\frac{3l(q,\frac{\bar{\tau}}{2})}{(\frac{\bar{\tau}}{2})}\\
&  \leq
\frac{6}{\bar{\tau}}\left(\sqrt{\frac{n}{2}}+\sqrt{\frac{3A}{2}}\right)^2,\nonumber
\end{align}  %\eqno (6.2.8)
for $q\in
B_{t_0-\frac{\bar{\tau}}{2}}(q(\frac{\bar{\tau}}{2}),\sqrt{A\bar{\tau}})$.
Recall that the Li-Yau-Hamilton inequality implies that the scalar
curvature of the ancient solution is pointwise nondecreasing in
time. Thus we know from (6.2.8) that \begin{equation} \tau R(q,t_0-\tau)\leq
6A\left(\sqrt{\frac{n}{2}}+\sqrt{\frac{3A}{2}}\right)^2 %\eqno(6.2.9)
\end{equation} whenever $\frac{1}{2}\bar{\tau}\leq \tau \leq A\bar{\tau}$ and
$d_{t_0-\frac{\bar{\tau}}{2}}^2 (q,q(\frac{\bar{\tau}}{2}))\leq
A\bar{\tau}.$

By (6.2.2) and (6.2.3) we have
$$
\frac{\partial l}{\partial \tau}+\frac{1}{2}|\nabla
l|^2=-\frac{l}{2\tau}+\frac{R}{2}.
$$
This together with (6.2.9) implies that
\begin{align*}
\frac{\partial l}{\partial \tau} &\leq
-\frac{l}{2\tau}+\frac{3A}{\tau}
\left(\sqrt{\frac{n}{2}}+\sqrt{\frac{3A}{2}}\right)^2 \\
\intertext{i.e.,} \frac{\partial}{\partial\tau}(\sqrt{\tau}l)
&\leq\frac{3A}{\sqrt{\tau}}\left(\sqrt{\frac{n}{2}}+\sqrt{\frac{3A}{2}}\right)^2
\end{align*}
whenever $\frac{1}{2}\bar{\tau}\leq \tau \leq A\bar{\tau}$ and
$d_{t_0-\frac{\bar{\tau}}{2}}^2
\left(q,q\left(\frac{\bar{\tau}}{2}\right)\right)\leq A\bar{\tau}.$ Hence by
integrating this differential inequality, we obtain
$$
\sqrt{\tau}l(q,\tau)-\sqrt{\frac{\bar{\tau}}{2}}l\left(q,\frac{\bar{\tau}}{2}\right)
\leq 6A\left(\sqrt{\frac{n}{2}}+\sqrt{\frac{3A}{2}}\right)^2\sqrt{\tau}
$$
and then by (6.2.7),
\begin{align}
l(q,\tau)&  \leq
l\left(q,\frac{\bar{\tau}}{2}\right)+6A\left(\sqrt{\frac{n}{2}}+\sqrt{\frac{3A}{2}}\right)^2\\
&  \leq 7A\left(\sqrt{\frac{n}{2}}+\sqrt{\frac{3A}{2}}\right)^2 \nonumber
\end{align}  %\eqno (6.2.10)
whenever $\frac{1}{2}\bar{\tau}\leq \tau \leq A\bar{\tau}$ and
$d_{t_0-\frac{\bar{\tau}}{2}}^2 (q,q(\frac{\bar{\tau}}{2}))\leq
A\bar{\tau}.$  So we have proved claim (6.2.6).
\end{proof}

\begin{lemma}[finite parabolic curvature point selection]
  \label{lem:cz-finite-parabolic-curvature-point-selection}
  \dcref{Assertion (6.3.4)}
  \group{Curvature Estimates from Volume Growth}
  \source{cao-zhu}
We first claim that there exists a point $(\bar{x},\bar{t})$ with
$-t_0<\bar{t}\leq 0$ and $d_{\bar{t}}(\bar{x},x_0)<\frac{1}{3}$
such that
$$
Q=|Rm(\bar{x},\bar{t})|>C+ B(\bar{t}+t_0)^{-1},
$$
and \begin{equation}
|Rm(x,t)|\leq 4Q %\eqno (6.3.1)? ---> 6.3.4     WRONG number
\end{equation} wherever $\bar{t}-AQ^{-1}\leq t\leq \bar{t},\ \ d_t(x,x_0)\leq
d_{\bar{t}}(\bar{x},x_0) +(AQ^{-1})^{\frac{1}{2}}.$
\end{lemma}
\begin{proof}
We will construct such $(\bar{x},\bar{t})$ as a limit of a finite
sequence of points.  Take an arbitrary $(x_1,t_1)$ such that
$$
d_{t_1}(x_1,x_0)\leq \frac{1}{4},\ \ \ -t_0<t_1\leq 0 \; \mbox{
and } \;|Rm(x_1,t_1)|>C+B(t_1+t_0)^{-1}.
$$
Such a point clearly exists by our assumption. Assume we have
already constructed $(x_k,t_k).$ If we cannot take the point
$(x_k,t_k)$ to be the desired point $(\bar{x},\bar{t})$, then
there exists a point $(x_{k+1},t_{k+1})$ such that
$$
t_k-A|Rm(x_k,t_k)|^{-1}\leq t_{k+1}\leq t_k,
$$
and
$$
d_{t_{k+1}}(x_{k+1},x_0) \leq
d_{t_{k}}(x_{k},x_0)+(A|Rm(x_k,t_k)|^{-1})^{\frac{1}{2}},
$$
but
$$
|Rm(x_{k+1},t_{k+1})|>4|Rm(x_k,t_k)|.$$

\vskip 0.1cm\noindent It then follows that
\begin{align*}
d_{t_{k+1}}(x_{k+1},x_0)& \leq
d_{t_{1}}(x_{1},x_0)+A^{\frac{1}{2}}
\left(\sum\limits_{i=1}^k|Rm(x_i,t_i)|^{-\frac{1}{2}}\right)\\
&  \leq \frac{1}{4}+A^{\frac{1}{2}}
\left(\sum\limits_{i=1}^k2^{-(i-1)}|Rm(x_1,t_1)|^{-\frac{1}{2}}\right)\\
&  \leq \frac{1}{4}+2(AC^{-1})^{\frac{1}{2}} \\
&  < \frac{1}{3},
\end{align*}
\begin{displaymath}
\begin{split}
t_{k+1}-(-t_0)&  =\sum_{i=1}^k(t_{i+1}-t_i)+(t_1-(-t_0))\\[1mm]
&  \geq-\sum_{i=1}^kA|Rm(x_i,t_i)|^{-1}+(t_1-(-t_0))\\[1mm]
&  \geq-A\sum_{i=1}^k4^{-(i-1)}|Rm(x_1,t_1)|^{-1}+(t_1-(-t_0))\\[1mm]
&  \geq-\frac{2A}{B}(t_1+t_0)+(t_1+t_0)\\[1mm]
&  \geq\frac{1}{2}(t_1+t_0),
\end{split}
\end{displaymath}
and
\begin{align*}
|Rm(x_{k+1},t_{k+1})|&  > 4^k|Rm(x_1,t_1)| \\
&  \geq 4^kC\rightarrow +\infty\; \mbox{ as } \;
k\rightarrow+\infty.
\end{align*}
Since the solution is smooth, the sequence $\{(x_k, t_k)\}$ is
finite and its last element fits. Thus we have proved  assertion
\end{proof}
